\documentclass[12pt, a4paper,journal, onecolumn, compsoc]{IEEEtran}
\IEEEoverridecommandlockouts
% The preceding line is only needed to identify funding in the first footnote. If that is unneeded, please comment it out.
\usepackage{cite}
\usepackage{amsmath,amssymb,amsfonts}
\usepackage{algorithmic}
\usepackage{graphicx}
\usepackage{textcomp}
\usepackage{xcolor}
\usepackage{fancyhdr}
\usepackage{mathtools}
\usepackage{lastpage}
\usepackage[bookmarks=false,hypertexnames=false]{hyperref}
\usepackage{cancel}
\usepackage{bm}
\usepackage{float}
\usepackage[RPvoltages]{circuitikz}
\usepackage{thmtools}
\usepackage{thm-restate}
\usepackage[titletoc,title]{appendix}
\hypersetup{
    colorlinks=true,
    linkcolor=blue,
    filecolor=magenta,      
    urlcolor=cyan,
    citecolor=blue,
}
\pagestyle{fancy}


\fancyhf{}
\fancyhead[C]{\small UNCuyo -- Ingeniería Mecatrónica / Mendoza -- Argentina. Espacio Curricular 311:\\ AUTOMÁTICA Y MÁQUINAS ELÉCTRICAS / PROYECTO GLOBAL INTEGRADOR.\\ Año: 2025. Alumnos: Calderón, Joaquín - Castel, Francisco. Fecha: 19 de Febrero 2026.}
\fancyfoot[C]{\small Pág. \thepage \hspace{1pt} de \pageref{LastPage}}
\renewcommand{\headrulewidth}{0.4pt}
\renewcommand{\footrulewidth}{0.4pt}

\setlength{\headheight}{45pt}
\def\BibTeX{{\rm B\kern-.05em{\sc i\kern-.025em b}\kern-.08em
    T\kern-.1667em\lower.7ex\hbox{E}\kern-.125emX}}
\makeatletter
\let\ps@plain\ps@fancy
\makeatother

% --- AJUSTE DE ESPACIADO DE TEXTO ---
\setlength{\parskip}{1em} % Añade espacio entre párrafos (y por ende, entre secciones)
\renewcommand{\baselinestretch}{1.15} % Aumenta un 15% el espacio entre líneas (más legible)

\begin{document}

% --- AJUSTE DE ESPACIADO DE ECUACIONES (GLOBAL) ---
\setlength{\abovedisplayskip}{12pt plus 3pt minus 3pt}      % Espacio antes de ecuación
\setlength{\belowdisplayskip}{12pt plus 3pt minus 3pt}      % Espacio después de ecuación
\setlength{\abovedisplayshortskip}{10pt plus 3pt minus 3pt} % Espacio corto antes
\setlength{\belowdisplayshortskip}{10pt plus 3pt minus 3pt} % Espacio corto después
\hbadness=10000
\hfuzz=10pt
\emergencystretch=3em

% =====================================================================
%  PORTADA
% =====================================================================
\begin{titlepage}
\thispagestyle{empty}
\centering

\vspace*{1cm}

{\large \textbf{Universidad Nacional de Cuyo}}\\[4pt]
{\large Facultad de Ingeniería}\\[4pt]
{\normalsize Ingeniería Mecatrónica}

\vspace{0.8cm}
\rule{\textwidth}{1.5pt}

\vspace{1.2cm}

{\huge \textbf{Proyecto Global Integrador}}\\[12pt]
{\LARGE Control de Accionamiento de CA con\\Motor Sincrónico de Imanes Permanentes}

\vspace{1.2cm}
\rule{\textwidth}{1.5pt}

\vspace{1.5cm}

{\large \textbf{Espacio Curricular 311}}\\[6pt]
{\Large Automática y Máquinas Eléctricas}

\vspace{2cm}

\begin{tabular}{r l}
    \textbf{Alumnos:} & Francisco Castel \\
                      & \texttt{castel.francisco@uncuyo.edu.ar} \\[8pt]
                      & Joaquín Calderón \\
                      & \texttt{calderon.joaquin@uncuyo.edu.ar} \\[12pt]
    \textbf{Año lectivo:} & 2025 \\[8pt]
    \textbf{Fecha:} & 19 de Febrero de 2026 \\[8pt]
    \textbf{Lugar:} & Mendoza, Argentina \\
\end{tabular}

\vfill

\rule{0.6\textwidth}{0.4pt}\\[6pt]
{\small Facultad de Ingeniería -- UNCuyo}

\end{titlepage}

% =====================================================================
%  CONTENIDO (comienza en página 2)
% =====================================================================
\setcounter{page}{1}
\thispagestyle{fancy}

\vspace*{-1cm}
\noindent\textbf{\large Abstract}
\vspace{4pt}

\noindent\textit{Modelado, simulación, diseño y análisis de desempeño de un Sistema de Control Automático de
Posición y Movimientos para un Accionamiento Eléctrico de 4 cuadrantes, compuesto por: máquina
eléctrica de corriente alterna (CA) trifásica sincrónica con excitación por imanes permanentes (PMSM),
alimentada por un inversor electrónico trifásico desde fuente de tensión continua (CC); reductor de
velocidad con engranajes planetarios conectado a la carga mecánica; retroalimentación con 1 sensor
de posición en el eje del motor más 3 sensores de corriente instantánea de fases en la salida del
inversor trifásico al estator de la máquina PMSM y 1 sensor de temperatura en el devanados de estator.}

\vspace{6pt}
\noindent\textbf{Keywords:} accionamiento CA, simulación, térmico, cascada, control, modelado.

\vspace{8pt}
\noindent\rule{\textwidth}{0.4pt}
\vspace{4pt}

\noindent\textbf{\large Resumen}
\vspace{4pt}

\noindent Este trabajo presenta el modelado, simulación y diseño de un sistema de control de posición y movimiento para un accionamiento eléctrico basado en un motor sincrónico de imanes permanentes (PMSM) trifásico acoplado mediante un reductor planetario ($r = 120$:1) a una carga mecánica tipo péndulo rígido de un grado de libertad. Se desarrolla el modelo no lineal completo de 6 estados en coordenadas rotóricas $qd0$ (incluyendo dinámica térmica del estator), y se obtiene un modelo LTI aumentado de 4 estados mediante linealización por retroalimentación de estado no lineal, empleando las leyes de control de desacoplamiento mínima ($v_{ds} = -L_q i_{qs} \omega_r$) y complementaria ($v_{qs} = v_{qs}^* + \omega_r L_d i_{ds}$). Se valida la equivalencia entre ambos modelos mediante simulación dinámica comparativa, verificando la concordancia de las respuestas ante secuencias de pulsos de tensión y torque de carga que exploran los cuatro cuadrantes de operación. Se analiza el efecto de condiciones iniciales $i_{ds}^r(0) \neq 0$ y de la inyección de tensión en eje $d$ (field forcing/weakening). Sobre el modelo LTI se diseña un controlador de movimiento en cascada con modulador de torque, cuyo desempeño se evalúa en simulación ante diversos escenarios de operación.

\vspace{4pt}
\noindent\rule{\textwidth}{0.4pt}

\section{Introducción}
El presente trabajo constituye el Proyecto Global Integrador del espacio curricular "Automática y Máquinas Eléctricas", cuyo objetivo principal es la integración de conocimientos y competencias fundamentales mediante su aplicación en un sistema mecatrónico.
El sistema bajo estudio consiste en un accionamiento eléctrico de cuatro cuadrantes compuesto por una máquina de corriente alterna trifásica sincrónica con excitación por imanes permanentes (PMSM), alimentada por un inversor electrónico. Este sistema acciona, a través de un reductor de velocidad, una carga mecánica correspondiente a una articulación de un brazo manipulador robótico elemental de un grado de libertad (péndulo rígido).
El alcance del proyecto abarca el modelado matemático, la simulación en Matlab/Simulink, y el diseño y análisis de desempeño de un sistema de control automático de posición y movimientos.

\section{Análisis, Modelado y Simulación del subsistema físico a Lazo Abierto}
Modelado de la carga mecánica, péndulo rígido de 1 g.d.l., fijo a una base en el sistema de referencia inercial, sometido a la fuerza de gravedad \(g\) con carga útil variable.

% 1. ECUACIONES DEL PÉNDULO
\begin{align}
    J_l\cdot \frac{d\omega_l}{dt} = T_q - b_l\cdot \omega_l - T_l(t)  \\
    T_l(t) = g\cdot k_l \cdot \sin(\theta_l(t)) + T_{ld}(t)
\end{align}

% 2. TEXTO EXPLICATIVO (Sin interrupciones)
Donde $\omega_l(t)$ y $\theta_l(t)$ representan la velocidad y posición angular de la articulación respectivamente; $T_q(t)$ es el torque impulsor entregado por el reductor y $T_{ld}(t)$ modela las perturbaciones externas variables. Los parámetros físicos constantes que caracterizan al péndulo y la transmisión se detallan con sus valores numéricos en la Tabla \ref{tab:parametros}.

% 3. CONTINUACIÓN LÓGICA (Inmediatamente después)
Es evidente que la relación cinemática viene dada por:
\[
    \frac{d\theta_l(t)}{dt} = \omega_l(t) \Leftrightarrow \theta_l(t) = \int_0^t{\omega_l(\xi)d\xi + \theta_l(0)}
\]

% 4. CÓDIGO DE LA TABLA (Al final, para que flote libremente)
\begin{table}[htbp]
    \caption{Parámetros Físicos y Especificaciones del Sistema}
    \centering
    \renewcommand{\arraystretch}{1.2} % Da un poco más de aire a las filas
    \begin{tabular}{|c|l|c|c|}
        \hline
        \textbf{Símbolo} & \textbf{Descripción} & \textbf{Valor} & \textbf{Unidad} \\
        \hline
        \multicolumn{4}{|c|}{\textbf{Subsistema Mecánico (Carga y Transmisión)}} \\
        \hline
        $r$ & Relación de transmisión & 120:1 & - \\
        $g$ & Aceleración de la gravedad & 9.80665 & $\text{m}/\text{s}^2$ \\
        $m$ & Masa del brazo manipulador & 1.0 & kg \\
        $l_{cm}$ & Longitud al centro de masa & 0.25 & m \\
        $J_{cm}$ & Inercia centro de masa brazo & 0.0208 & $\text{kg}\cdot\text{m}^2$ \\
        $m_l$ & Masa de carga útil (variable) & $0 - 1.5$ & kg \\
        $l_{l}$ & Longitud total al extremo & 0.50 & m \\
        $b_{l}$ & Fricción viscosa articulación & 0.1 & $\text{Nms}/\text{rad}$ \\
        \hline
        \multicolumn{4}{|c|}{\textbf{Subsistema Mecánico (Motor)}} \\
        \hline
        $J_m$ & Inercia rotor (motor + caja) & $14 \times 10^{-6}$ & $\text{kg}\cdot\text{m}^2$ \\
        $b_m$ & Fricción viscosa (motor + caja) & $15 \times 10^{-6}$ & $\text{Nms}/\text{rad}$ \\
        $J_{eq}$ & Inercia equiv. total (Calc.) & $1.978 \times 10^{-5}$ & $\text{kg}\cdot\text{m}^2$ \\
        $b_{eq}$ & Fricción equiv. total (Calc.) & $2.194 \times 10^{-5}$ & $\text{Nms}/\text{rad}$ \\
        $k_l$ & Coef. gravitacional de carga (Calc.) & $0.25$ & $\text{kg}\cdot\text{m}$ \\
        \hline
        \multicolumn{4}{|c|}{\textbf{Subsistema Eléctrico (PMSM)}} \\
        \hline
        $P_p$ & Pares de polos & 3 & - \\
        $\lambda'_m$ & Flujo magnético de imanes & 0.016 & Wb \\
        $L_d$ & Inductancia eje directo & 6.6 & mH \\
        $L_q$ & Inductancia eje cuadratura & 5.8 & mH \\
        $L_{ls}$ & Inductancia de dispersión & 0.8 & mH \\
        $R_{sREF}$ & Resistencia estator (@ $20^\circ$C) & 1.02 & $\Omega$ \\
        \hline
        \multicolumn{4}{|c|}{\textbf{Subsistema Térmico}} \\
        \hline
        $C_{ts}$ & Capacitancia térmica estator & 0.818 & $\text{J}/^\circ\text{C}$ \\
        $R_{ts-amb}$ & Resistencia térmica estator-ambiente & 146.7 & $^\circ\text{C}/\text{W}$ \\
        $\alpha_{Cu}$ & Coef. temp. del cobre & $3.9 \times 10^{-3}$ & $1/^\circ\text{C}$ \\
        $T_{amb}$ & Temperatura ambiente & 20 & $^\circ\text{C}$ \\
        \hline
        \multicolumn{4}{|c|}{\textbf{Límites de Operación (Especificaciones)}} \\
        \hline
        $n_{l,nom}$ & Velocidad nominal de carga & 60 & rpm \\
        $T_{q,nom}$ & Torque nominal de salida & 17.0 & Nm \\
        $T_{q,max}$ & Torque máximo (pico) & 45.0 & Nm \\
        \hline
    \end{tabular}
    \label{tab:parametros}
\end{table}

\subsection{Tren de transmisión}
Como se explicó anteriormente, la velocidad del motor no tendrá relación 1:1 con el péndulo, por lo que se implementa un \textbf{tren de transmisión planetario} como caja reductora reversible. Se asume lo siguiente:
\begin{itemize}
    \item Ausencia de \textit{backlash}.
    \item Ausencia de elasticidad torcional (acoplamiento rígido).
    \item Momento de inercia equivalente (parámetros concentrados).
    \item Pérdidas de energía por fricción interna consideradas en el eje de entrada junto al motor.
\end{itemize}
Por lo tanto, debido a las simplificaciones realizadas, es válido considerar a la carga mecánica y el motor eléctrico como una unidad en sí misma, lo que permite tener un subsistema mecánico de un grado de libertad, permitiendo escribir las siguientes ecuaciones.
\begin{align}
    \omega_l(t) = \frac{1}{r}\omega_m(t) \\
    T_q(t) = rT_d(t)
\end{align}
tal que:
\begin{itemize}
    \item r: Relación de velocidades del tren de transmisión.
    \item \(\omega_m(t)\): Velocidad angular del eje del motor.
    \item \(T_d(t)\): Torque de entrada al tren de transmisión.
    \item \(T_q(t)\): Torque de salida al tren de transmisión.
    
\end{itemize}

    \begin{figure}[htbp]
        \centering
        \includegraphics[width=0.6\linewidth]{imagenes/brazo_1gdl.png}
        \caption{Robot manipulador elemental de 1 g.d.l en plano vertical (péndulo rígido actuado)}
        \label{fig:brazo_1gdl}
    \end{figure}

\subsection{Máquina PMSM}
    Para el accionamiento, se utilizará una máquina eléctrica de CA trifásica sincrónica con excitación por imanes permanentes (\textbf{PMSM}) y estator simétrico y equilibrado, con bornes de fases \textit{abc} (coordenadas reales y "estacionarias) conectados en configuración estrella (Y) con neutro flotante \textbf{no} accesible.
    \subsubsection{Modelo matemático equivalente del subsistema mecánico}
    Habiendo expresado las ecuaciones para los distintos subsistemas de la máquina PMSM y del tren de transmisión, es posible conjugarlas en un modelo equivalente.
    Sustituyendo \(\omega_l(t)\) de (3) y \(T_q(t) = rT_d(t)\) de (4) en la Ec. (1) tenemos:
    \begin{equation}\label{EqAux1}
        \frac{J_l}{r}\frac{d\omega_m(t)}{dt} = rT_d(t)-\frac{b_l}{r}\omega_m(t)-T_l(t)
    \end{equation}
    Despejando \(T_d(t)\) y reemplazando \(T_l(t)\) por la Eq. 2:
    \begin{equation}\label{Eq1}
        T_d(t)=\frac{J_l}{r^2}\frac{d\omega_m}{dt}+\frac{b_l}{r^2}\omega_m+\frac{1}{r}\left(g\cdot k_l\cdot \sin\Big(\frac{\theta_m(t)}{r}\Big)+T_{ld}(t)\right)
    \end{equation}
    Por otra parte, la ecuación mecánica del rotor del motor es:
    \begin{equation}\label{EqJm}
        J_m\frac{d\omega_m}{dt}=T_m-b_m\omega_m-T_d(t)
    \end{equation}
    Reemplazando \(T_d(t)\) de \eqref{Eq1} en  \eqref{EqJm}:
    \begin{equation}
        J_{eq}\frac{d\omega_m}{dt}=T_m-b_m\omega_m-\frac{1}{r}\Big(g\cdot k_l \sin{\Big(\frac{\theta_m(t)}{r}\Big)} + T_{ld}(t)\Big)
    \end{equation}
    \begin{gather}
        J_{eq} = J_m +\frac{J_l}{r^2}\\
        b_{eq} = b_m + \frac{b_l}{r^2}\\
        \theta_l(t)=\frac{\theta_m(t)}{r}
    \end{gather}
    Podemos entonces, a partir de lo desarrollado, expresar el modelo en el espacio de estados.

    \textbf{Ecuaciones de estado:} Las ecuaciones de estado del subsistema mecánico completo quedan:
    \begin{equation}
        \begin{cases}
            \dot \theta_m(t) = \omega_m(t), \\ 
            \dot\omega_m(t)=-\frac{b_{eq}}{J_{eq}}\omega_m(t)-\frac{g\cdot k_l}{J_{eq}r}\cdot \sin{\Big(\frac{\theta_m(t)}{r}\Big)}+\frac{1}{J_{eq}}T_m(t)-\frac{1}{J_{eq}r}T_{ld}(t),\\
            y(t)=\theta_m(t)
        \end{cases}
        \label{eq: EqEstadoModeloMecanico}
    \end{equation}
    \textbf{Forma matricial}
    \begin{equation}
        \begin{cases}
            \begin{bmatrix}
                \dot\theta_m(t) \\\dot\omega_m(t)
            \end{bmatrix}
            = 
            \begin{bmatrix}
                \omega_m(t)\\
                -\frac{g\cdot k_l}{J_{eq}r}\cdot\sin{\big(\frac{\theta_m(t)}{r}\big)}-\frac{b_{eq}}{J_{eq}}\cdot\omega_m(t)
            \end{bmatrix}
            +
            \begin{bmatrix}
                0 & 0\\
                \frac{1}{J_{eq}} &-\frac{1}{J_{eq}r}
            \end{bmatrix}
            \begin{bmatrix}
                T_m(t)\\
                T_{ld}(t)
            \end{bmatrix}\\
            y(t) = \begin{bmatrix}
                1 & 0
            \end{bmatrix}
            \begin{bmatrix}
                \theta_m(t)\\
                \omega_m(t)
            \end{bmatrix}
        \end{cases}
    \end{equation}
    Notesé como se presenta una no linealidad en el primer termino lo que impide formular la matriz $A$ de forma convencional.
    \subsubsection{Modelo del subsistema electromagnético del motor}
    Modelo idealizado equivalente en coordenadas rotóricas \(qd0\), fijas al rotor, obtenido aplicando la \textit{Transformación de Park} al circuito de estator estacionarion (marco de referencia virtual eléctrico de rotor, sincrónico solo en régimen estacionario).
    \begin{gather}
        \frac{d\theta_r(t)}{dt} \equiv \omega_r(t) \\
        \theta_r(t)=\int_0^t{\omega_r(\xi)d\xi+\theta_r(0)} \\
        \theta_r(t) \equiv P_p\theta_m(t) \therefore \omega_r(t) = P_p\omega_m(t)
    \end{gather}
    \textbf{Torque electromagnético} = \textit{torque debido a imanes + torque de reluctancia}:
    \begin{equation}
            T_m(t) = \frac{3}{2} P_p \cdot \lambda'^r_m \cdot i^r_{qs}(t) + \frac{3}{2} P_p (L_d-L_q)\cdot i^r_{ds}(t)\cdot i^r_{qs}(t)
            \label{eq:torque_electromagnetico}
    \end{equation}
\textbf{Ecuaciones de tensiones de fase referidas al estator}\\
Dado que es necesario generar un campo magnético rotatorio sinusoidal para efectuar un movimiento en los motores PMSM, y que además, físicamente los devanados están separados \(120^\circ\) entre sí, es de esperarse que el inversor trifásico deba entregar las siguientes tensiones al motor (sistema equilibrado).
\begin{align}
v_{as}(t) &= V_m \cos\big(\theta_{ev}(t)\big) \\
v_{bs}(t) &= V_m \cos\left(\theta_{ev}(t) - \frac{2\pi}{3}\right) \\
v_{cs}(t) &= V_m \cos\left(\theta_{ev}(t) + \frac{2\pi}{3}\right)
\end{align}
Por lo que definiremos:
\[\bm{\mathrm{v}}_{abcs} = (v_{as},\ v_{bs},\ v_{cs})\]
\textbf{Ecuaciones de tensiones de fase referidas al rotor}
    \begin{align}
        K_s\textbf{v}_{abcs} = \textbf{v}^r_{qd0s}
    \end{align}
    Siendo \(K_s\) la matriz de transformación de Park.
    \begin{equation}\label{eq:vqs_rotor}
        v^r_{qs}=R_s(T_s^\circ(t))i^r_{qs}(t) + L_q\frac{d i^r_{qs}(t)}{dt}+[\lambda'^r_m + L_d i^r_{ds}(t)]\omega_r(t)
    \end{equation}
    \begin{equation}\label{eq:vds_rotor}
        v_{ds}^r=R_s(T_s^\circ(t))i^r_{ds}(t)+L_d\frac{d i^r_{ds}(t)}{dt} - L_qi^r_{qs}(t)\omega_r(t)
    \end{equation}
    \begin{equation}\label{eq:v0s_rotor}
        v^r_{0s}=R_s(T_s^\circ(t))i^r_{0s}(t)+L_{ls}\frac{d i^r_{0s}(t)}{dt}
    \end{equation}

    \subsubsection{Modelo del subsistema térmico}
    El calentamiento del estator se modela mediante un circuito térmico concentrado de primer orden, con una capacitancia térmica \(C_{ts}\) asociada al almacenamiento de energía y una resistencia térmica \(R_{ts\text{-}amb}\) asociada a la transferencia de calor hacia el ambiente. La potencia que excita este subsistema proviene de las pérdidas Joule del bobinado estatórico.

    La resistencia del estator depende de la temperatura según:
    \begin{equation}
        R_s\big(T_s^\circ(t)\big)=R_{sREF}\Big[1+\alpha_{Cu}\big(T_s^\circ(t)-T_{REF}^\circ\big)\Big], \qquad T_{REF}^\circ=20^\circ\text{C}
    \end{equation}

    Bajo las hipótesis adoptadas, el balance de potencia térmica queda:
    \begin{equation}
        C_{ts}\frac{dT_s^\circ(t)}{dt}=P_{Cu}(t)-\frac{T_s^\circ(t)-T_{amb}^\circ(t)}{R_{ts\text{-}amb}}
    \end{equation}
    donde
    \begin{equation}
        P_{Cu}(t)=\frac{3}{2}R_s\big(T_s^\circ(t)\big)\Big(i_{qs}^{r\,2}(t)+i_{ds}^{r\,2}(t)+2i_{0s}^{r\,2}(t)\Big)
    \end{equation}

    Por lo tanto, la ecuación de estado del subsistema térmico es:
    \begin{equation}
        \frac{dT_s^\circ(t)}{dt}=\frac{1}{C_{ts}}\left[\frac{3}{2}R_s\big(T_s^\circ(t)\big)\Big(i_{qs}^{r\,2}(t)+i_{ds}^{r\,2}(t)+2i_{0s}^{r\,2}(t)\Big)-\frac{T_s^\circ(t)-T_{amb}^\circ(t)}{R_{ts\text{-}amb}}\right]
    \end{equation}

    \subsubsection{Modelo no lineal global del accionamiento}
    Una vez definidos el subsistema mecánico equivalente, el subsistema electromagnético en coordenadas \(qd0\) y el subsistema térmico, puede escribirse el modelo no lineal global del accionamiento. Tomando como vector de estado
    \[
        \mathbf{x}_{NL}(t)=\begin{bmatrix}
            \theta_m(t) & \omega_m(t) & i_{qs}^r(t) & i_{ds}^r(t) & i_{0s}^r(t) & T_s^\circ(t)
        \end{bmatrix}^T
    \]
    y como entradas
    \[
        \mathbf{u}_{NL}(t)=\begin{bmatrix}
            v_{qs}^r(t) & v_{ds}^r(t) & v_{0s}^r(t) & T_{ld}(t) & T_{amb}^\circ(t)
        \end{bmatrix}^T
    \]
    el sistema queda descrito por el siguiente conjunto de ecuaciones de estado no lineales:
    \begin{equation}
    \left\{
        \begin{aligned}
            \dot{\theta}_m(t) &= \omega_m(t) \\[6pt]
            \dot{\omega}_m(t) &= -\frac{b_{eq}}{J_{eq}}\omega_m(t)-\frac{gk_l}{J_{eq}r}\sin\left(\frac{\theta_m(t)}{r}\right)+\frac{1}{J_{eq}}T_m(t)-\frac{1}{J_{eq}r}T_{ld}(t) \\[6pt]
            \dot{i}_{qs}^r(t) &= \frac{1}{L_q}\Big[v_{qs}^r(t)-R_s\big(T_s^\circ(t)\big)i_{qs}^r(t)-\big(\lambda'_m+L_d i_{ds}^r(t)\big)\omega_r(t)\Big] \\[6pt]
            \dot{i}_{ds}^r(t) &= \frac{1}{L_d}\Big[v_{ds}^r(t)-R_s\big(T_s^\circ(t)\big)i_{ds}^r(t)+L_q i_{qs}^r(t)\omega_r(t)\Big] \\[6pt]
            \dot{i}_{0s}^r(t) &= \frac{1}{L_{ls}}\Big[v_{0s}^r(t)-R_s\big(T_s^\circ(t)\big)i_{0s}^r(t)\Big] \\[6pt]
            \dot{T}_s^\circ(t) &= \frac{1}{C_{ts}}\left[\frac{3}{2}R_s\big(T_s^\circ(t)\big)\Big(i_{qs}^{r\,2}(t)+i_{ds}^{r\,2}(t)+2i_{0s}^{r\,2}(t)\Big)-\frac{T_s^\circ(t)-T_{amb}^\circ(t)}{R_{ts\text{-}amb}}\right]
        \end{aligned}
    \right.
    \end{equation}
    con
    \begin{equation}
        T_m(t)=\frac{3}{2} P_p \lambda'_m i_{qs}^r(t)+\frac{3}{2}P_p(L_d-L_q)i_{ds}^r(t)i_{qs}^r(t),
        \qquad
        \omega_r(t)=P_p\omega_m(t)
    \end{equation}

    Este modelo de 6 estados será la planta de referencia para la implementación en Simulink y para las linealizaciones posteriores. Además, deja explícito que las no linealidades relevantes del sistema provienen del término gravitacional \(\sin(\theta_m/r)\), del acoplamiento electromagnético entre ejes \(q\) y \(d\), y de la dependencia térmica de \(R_s\big(T_s^\circ\big)\).

    % ==========================================================================
    % DEMOSTRACIÓN: La ecuación de secuencia cero no influye con neutro flotante
    % ==========================================================================

    \subsubsection{Irrelevancia de la ecuación de secuencia cero en el modelo del accionamiento}

    La ecuación de secuencia cero del modelo electromagnético en coordenadas rotóricas queda dada por:
    \begin{equation}
        v_{0s}^r(t)=R_s\!\big(T_s^\circ(t)\big)\,i_{0s}^r(t)+L_{ls}\frac{d i_{0s}^r(t)}{dt}
        \label{eq:v0s_rotor_final}
    \end{equation}

    De acuerdo con la consigna, se demostrará que esta ecuación no influye en la dinámica del sistema bajo dos situaciones distintas:

    \begin{enumerate}
        \item sistema trifásico equilibrado con circuito electromagnético de estator simétrico y retorno de neutro;
        \item sistema general no necesariamente simétrico ni equilibrado, pero con centro de estrella flotante (no accesible).
    \end{enumerate}

    % Figura: Conexión estrella con neutro flotante — LCK en nodo n
    \begin{figure}[H]
    \centering
    \begin{circuitikz}[scale=0.95, transform shape, american currents, european resistors]

        % Nodo neutro
        \coordinate (N) at (0,0);
        \filldraw (N) circle (2pt);
        \node[below left=2pt] at (N) {$n$};

        %========================
        % FASE A
        %========================
        \draw
        (N) -- (0,1.2)
        to[L, l=$L_a$, i>^=$i_a$] (0,3.0)
        to[R, l=$R_s$] (0,4.6)
        -- (0,5.2);
        \node[above] at (0,5.2) {$a$};

        % Tensión fase-neutro
        \draw[blue,->,thick] (-1.1,0.15) -- (-1.1,5.05);
        \node[blue,left] at (-1.1,2.6) {$v_{an}$};

        %========================
        % FASE B
        %========================
        \draw
        (N) -- (-1.15,-0.7)
        to[L, l_=$L_b$, i>_=$i_b$] (-2.8,-1.75)
        to[R, l_=$R_s$] (-4.3,-2.65)
        -- (-4.9,-3.0);
        \node[left] at (-4.9,-3.0) {$b$};

        %========================
        % FASE C
        %========================
        \draw
        (N) -- (1.15,-0.7)
        to[L, l=$L_c$, i>^=$i_c$] (2.8,-1.75)
        to[R, l=$R_s$] (4.3,-2.65)
        -- (4.9,-3.0);
        \node[right] at (4.9,-3.0) {$c$};

        %========================
        % Neutro flotante
        %========================
        \draw[densely dotted, gray, thick] (N) -- (0,-1.3);
        \node[gray, below] at (0,-1.35) {\footnotesize neutro flotante};

    \end{circuitikz}
    \caption{Estator trifásico conectado en estrella con neutro flotante. Al no existir conductor de retorno externo, la Ley de Corrientes de Kirchhoff aplicada en el nodo neutro impone \(i_a+i_b+i_c=0\), por lo que la componente homopolar de corriente resulta nula: \(i_{0s}(t)\equiv 0\).}
    \label{fig:estrella_neutro_flotante}
    \end{figure}

    \textbf{Definición de la componente homopolar mediante la Transformación de Park.}
    La transformación \(abc \rightarrow qd0\) separa las variables trifásicas en una componente en cuadratura, una componente directa y una componente homopolar. En particular, la componente de secuencia cero de una magnitud genérica \(x\in\{v,i\}\) se define como:
    \begin{equation}
        x_{0s}(t)=\frac{1}{3}\Big(x_a(t)+x_b(t)+x_c(t)\Big)
        \label{eq:def_x0}
    \end{equation}
    Por lo tanto:
    \begin{equation}
        i_{0s}(t)=\frac{1}{3}\Big(i_a(t)+i_b(t)+i_c(t)\Big),
        \qquad
        v_{0s}(t)=\frac{1}{3}\Big(v_a(t)+v_b(t)+v_c(t)\Big)
        \label{eq:def_i0_v0}
    \end{equation}

    \textbf{Caso 1: sistema trifásico equilibrado y estator simétrico, con retorno de neutro.}
    Para un sistema trifásico equilibrado se cumple, en todo instante:
    \begin{equation}
        i_a(t)+i_b(t)+i_c(t)=0
        \label{eq:suma_i_eq}
    \end{equation}
    y, de manera análoga, para tensiones equilibradas:
    \begin{equation}
        v_a(t)+v_b(t)+v_c(t)=0
        \label{eq:suma_v_eq}
    \end{equation}

    Sustituyendo \eqref{eq:suma_i_eq} y \eqref{eq:suma_v_eq} en \eqref{eq:def_i0_v0}, resulta:
    \begin{equation}
        i_{0s}(t)\equiv 0,
        \qquad
        v_{0s}(t)\equiv 0
    \end{equation}

    En consecuencia, la ecuación \eqref{eq:v0s_rotor_final} se reduce a:
    \begin{equation}
        0=R_s\!\big(T_s^\circ(t)\big)\cdot 0 + L_{ls}\frac{d(0)}{dt}
    \end{equation}
    es decir, a una identidad trivial. Por lo tanto, en el caso equilibrado y simétrico, la ecuación de secuencia cero no aporta dinámica al modelo.

    \textbf{Caso 2: sistema general con centro de estrella flotante (no accesible).}
    Considérese ahora el caso general en el que no se supone ni simetría perfecta ni equilibrio de las excitaciones, pero el punto neutro del estator conectado en estrella es flotante, es decir, no existe conductor de retorno desde el centro de estrella hacia una referencia externa.

    \textbf{Demostración mediante Ley de Corrientes de Kirchhoff (LCK).}
    En el nodo neutro \(n\) confluyen únicamente las tres corrientes de fase. Como no existe un cuarto conductor de retorno, la LCK impone:
    \begin{equation}
        i_a(t)+i_b(t)+i_c(t)=0
        \label{eq:LCK_neutro_flotante}
    \end{equation}

    Esta restricción es puramente topológica y no depende de que el sistema sea o no equilibrado. Por lo tanto, usando nuevamente la definición \eqref{eq:def_i0_v0}:
    \begin{equation}
        i_{0s}(t)=\frac{1}{3}\Big(i_a(t)+i_b(t)+i_c(t)\Big)=0
    \end{equation}
    o sea,
    \begin{equation}
        \boxed{i_{0s}(t)\equiv 0}
        \label{eq:i0_neutro_flotante}
    \end{equation}

    \textbf{Interpretación mediante la Transformación de Park:}
    La componente homopolar obtenida por Park representa la parte común a las tres fases. Si la suma instantánea de corrientes de fase es nula, entonces la proyección sobre el eje \(0\) también es nula. Es decir, el subespacio homopolar no queda excitado:
    \begin{equation}
        i_{0s}(t)\equiv 0
    \end{equation}

    Sustituyendo este resultado en \eqref{eq:v0s_rotor_final}, se obtiene:
    \begin{equation}
        v_{0s}^r(t)=R_s\!\big(T_s^\circ(t)\big)\cdot 0 + L_{ls}\frac{d(0)}{dt}=0
    \end{equation}
    de modo que la ecuación de secuencia cero vuelve a quedar satisfecha de forma trivial, sin aportar dinámica propia al sistema.

    En consecuencia, la componente homopolar no aporta dinámica relevante para el análisis electromecánico del accionamiento bajo las hipótesis adoptadas. Este resultado justifica que, en los desarrollos posteriores, el eje \(0\) pueda omitirse cuando se trabaje con el modelo reducido en coordenadas \(qd\).

% ==========================================================================
% SECCIÓN II-E: VALIDACIÓN POR SIMULACIÓN (LAZO ABIERTO)
% ==========================================================================

\subsection{Implementación y Validación en Entorno de Simulación}
Se implementó el modelo matemático completo en el entorno MATLAB/Simulink, estructurando el sistema mediante bloques funcionales interconectados que representan las ecuaciones diferenciales deducidas previamente.

La arquitectura de la simulación se diseñó para validar tanto la dinámica de la planta como la coherencia de las transformaciones de coordenadas. El flujo de señales se configura de la siguiente manera:
\begin{enumerate}
    \item \textbf{Generación de Consigna:} Se establecen tensiones de referencia constantes en el marco rotórico ($v_q^*=5$ V, $v_d^*=0$).
    \item \textbf{Transformaciones de Coordenadas:} Las tensiones $dq$ se convierten al marco trifásico $abc$ mediante la \textit{Transformada Inversa de Park} y se reconvierten inmediatamente al marco $dq$ mediante la \textit{Transformada de Park}. Esta estructura en bucle ("Back-to-Back") permite verificar la integridad matemática de los algoritmos antes de ingresar a la planta.
    \item \textbf{Planta Física:} Las señales resultantes alimentan al \textbf{Subsistema Electromecánico} y al \textbf{Subsistema Mecánico}, interactuando simultáneamente con el \textbf{Subsistema Térmico}.
\end{enumerate}

% FIGURA DIAGRAMA SIMULINK
\begin{figure}[H]
    \centering
    \includegraphics[width=1.0\linewidth]{imagenes/simulink_diagrama_completo.png}
    \caption{Diagrama de bloques implementado en Simulink. Se destaca la estructura de validación con transformadas directa e inversa y la interconexión de los subsistemas.}
    \label{fig:simulink_top}
\end{figure}

\subsubsection{Resultados del Ensayo a Lazo Abierto}
Se realizó un ensayo de excitación al escalón con $v_q = 5$ V y carga nula ($T_{ld}=0$), partiendo de condiciones iniciales de reposo y equilibrio térmico ($T_s(0)=20^\circ$C).

\textbf{Análisis Electromecánico:}
La Fig. \ref{fig:res_elec_mec} vincula la dinámica eléctrica con la mecánica:
\begin{itemize}
    \item \textbf{Superior ($i_q$):} Al aplicar la tensión, la corriente aumenta rápidamente ($\tau_e = L_q/R_s$). Sin embargo, a medida que el motor gana velocidad, la aparición de la fuerza contraelectromotriz ($e = \lambda_r \omega_r$) se opone a la tensión aplicada, provocando la reducción de la corriente observada a partir de $t \approx 0.1$ s.
    \item \textbf{Inferior ($\omega_m$):} El torque generado por la corriente acelera el rotor venciendo la inercia $J_{eq}$. La velocidad se estabiliza cuando el torque motor se iguala con el torque de fricción viscosa.
\end{itemize}

% FIGURA COMBINADA CORRIENTE/VELOCIDAD (CLAVADA AQUÍ CON [H])
\begin{figure}[H]
    \centering
    \includegraphics[width=0.9\linewidth]{imagenes/resultado_corriente.png} \\[0.1cm]
    \includegraphics[width=0.9\linewidth]{imagenes/resultado_velocidad.png}
    \caption{Respuesta dinámica al escalón de tensión $v_q=5$V. \textbf{Arriba:} Corriente de estator $i_q$ [A], mostrando el transitorio inductivo y el efecto de la Back-EMF. \textbf{Abajo:} Velocidad mecánica del rotor $\omega_m$ [rad/s].}
    \label{fig:res_elec_mec}
\end{figure}

\textbf{Análisis Térmico:}
La Fig. \ref{fig:res_term} muestra la evolución de la temperatura del estator. Se observa un inicio correcto en $T_{amb}=20^\circ$C.
Es notable que la pendiente de calentamiento no es constante: es máxima al inicio (cuando la corriente es alta y el rotor está detenido) y disminuye cuando el motor acelera y la corriente baja debido a la fuerza contraelectromotriz. Este comportamiento valida el correcto acoplamiento entre las variables eléctricas ($i^2 R$) y la dinámica mecánica en el modelo térmico.

% FIGURA TÉRMICA (CLAVADA AQUÍ CON [H])
\begin{figure}[H]
    \centering
    \includegraphics[width=0.9\linewidth]{imagenes/resultado_temperatura.png}
    \caption{Evolución de la temperatura del estator $T_s$. Nótese el cambio de pendiente provocado por la variación de la corriente durante el arranque.}
    \label{fig:res_term}
\end{figure}

% ESTE COMANDO ASEGURA QUE LAS FOTOS TERMINEN AQUÍ Y NO PASEN A LA MATEMÁTICA
\clearpage



%-------------- LPV -----------------
\subsection{Linealización mediante el metodo de \textit{pequeñas señales}: Modelo LPV con \(i^r_{ds}\neq0\)}
A partir del modelo NL anteriormente presentado, es posible a través de la Linealización Jacobiana, mediante aproximación con Serie de Taylor truncada en 1er orden en un punto genérico de operación, obtener un modelo lineal mas simple de analizar y controlar.

Considerando que el actual modelo es de la forma \[\bm{\dot x} =\bm{f}(\bm{x},u)\]
Donde lo primero es encontrar los puntos de equilibrio, del sistema, es decir los valores \(\bm{x}_0 ,\)\(u_0 ,\) tal que \(\bm{\dot x_0} = \bm{f}(\bm{x}_0,u_0)\), y definimos \(x=x_0 +\delta x\) y \(u=u_0+\delta u\). Luego, expandimos la ecuación no lineal en términos de las perturbaciones respecto a estos valores de equilibrio, lo que resulta en: 
\begin{equation*}
    \dot x_0 + \delta \dot x \cong \bm{f}(\bm{x}_0,u_0) + A \delta x+B\delta u
\end{equation*}
Donde \(A\) y \(B\) son la mejor aproximación lineal de la función \(\bm{f}(x,u) \) en \(x_0, u_0\), calculándose: 
\begin{equation*}
    \bm{A}_0 = \left[ \frac{\partial \bm{f}}{\partial \bm{x}} \right]_{\bm{x}_o, u_o} \quad \text{y} \quad \bm{B}_0 = \left[ \frac{\partial \bm{f}}{\partial u} \right]_{\bm{x}_o, u_o}
\end{equation*}
Restando la solución de equilibrio y dejando la ecuación solo en variables incrementales nos queda:
\begin{equation*}
    \delta\dot x=A_0\delta \bm{x}+ B_0\delta u
\end{equation*}
Es decir, una ecuación diferencial lineal alrededor del punto de equilibrio cuasi-estático.
\subsubsection{Implementación en caso genérico de operación, \(\omega=cte\)}
Utilizando entonces las ecuaciones vistas podemos comenzar a linealizar nuestro sistema mediante el método de pequeñas señales.
Primero, encontramos nuestro espacio de operación global no lineal, que es el espacio donde se resolvera la primer parte de la ecuación expandida, $\bm{f}(\bm{x}_0,u_0)=0$

\begin{equation}
\left\{
    \begin{aligned}
        \frac{d\theta_{mO}}{dt}(t) &= \omega_{mO}(t) = \text{cte} \\[10pt]
        %
        \frac{d\omega_{mO}}{dt}(t) &= \frac{1}{J_{eq}} \left[ -b_{eq}\omega_{mO}(t) + \frac{3}{2}P_p \left[ \lambda_m' + (L_d - L_q)i_{dsO}^r(t) \right] i_{qsO}^r(t) - g \cdot k_l \sin\left(\frac{\theta_{mO}}{r}\right) - \frac{1}{r}T_{ld}(t) \right] = 0 \\[10pt]
        %
        \frac{di_{qsO}^r}{dt}(t) &= \frac{1}{L_q} \left[ v_{qsO}^r(t) - R_s(T_{sO}^\circ(t))i_{qsO}^r(t) - (\lambda_m' + L_d i_{dsO}^r(t))P_p \omega_{mO}(t) \right] = 0 \\[10pt]
        %
        \frac{di_{dsO}^r}{dt}(t) &= \frac{1}{L_d} \left[ v_{dsO}^r(t) - R_s(T_{sO}^\circ(t))i_{dsO}^r(t) + L_q i_{qsO}^r(t)P_p \omega_{mO}(t) \right] = 0 \\[10pt]
        %
        \frac{di_{0sO}^r}{dt}(t) &= \frac{1}{L_{ls}} \left( v_{0sO}^r(t) - R_s(T_{sO}^\circ(t))i_{0sO}^r(t) \right) = 0 \\[10pt]
        %
        \frac{dT_{sO}^\circ}{dt}(t) &= \frac{1}{C_{ts}} \left[ \frac{3}{2}R_s(t) \left( i_{qsO}^r(t)^2 + i_{dsO}^r(t)^2 + 2i_{0sO}^r(t)^2 \right) - \frac{1}{R_{ts-amb}} (T_{sO}^\circ(t) - T_{ambO}^\circ(t)) \right] = 0
    \end{aligned}
\right.
\end{equation}
Luego la parte lineal dinámica, que es de principal interés en este análisis.
\begin{equation*}
    \Delta \dot{x}(t) = \underbrace{\left. \frac{\partial f}{\partial x} \right|_0}_{A_o} (t) \cdot \Delta x(t) + \underbrace{\left. \frac{\partial f}{\partial u} \right|_0}_{B_o} (t) \cdot \Delta u(t); \quad \Delta x(0) = 0
\end{equation*}
\begin{equation*}
\left\{
    \begin{aligned}
        \Delta \dot{\theta}_m(t) &= \Delta \omega_m(t) \\[10pt]
        %
        \Delta \dot{\omega}_m(t) &= \frac{1}{J_{eq}} \left[ \frac{3}{2} P_p \left\{ [\lambda_m' + (L_d - L_q) \cdot i_{dsO}^r] \cdot \Delta i_{qs}^r(t) + (L_d - L_q) \cdot i_{qsO}^r \cdot \Delta i_{ds}^r(t) \right\} - b_{eq} \cdot \Delta \omega_m(t) \right. \\
        &\quad \left. - \frac{g \cdot k_l}{r^2} \cdot \cos\left(\frac{\theta_{mO}}{r}\right) \cdot \Delta \theta_m(t) - \frac{\Delta T_{ld}(t)}{r} \right] \\[10pt]
        %
        \Delta \dot{i}_{qs}^r(t) &= \frac{1}{L_q} \left[ \Delta v_{qs}^r(t) - R_{sO} \cdot \Delta i_{qs}^r(t) - R_{s,REF} \alpha_{cu} \Delta T_s^\circ(t) \cdot i_{qsO}^r - \{\lambda_m' + L_d \cdot i_{dsO}^r\} P_p \cdot \Delta \omega_m(t) \right. \\
        &\quad \left. - L_d P_p \omega_{mO} \cdot \Delta i_{ds}^r(t) \right] \\[10pt]
        %
        \Delta \dot{i}_{ds}^r(t) &= \frac{1}{L_d} \left[ \Delta v_{ds}^r(t) - R_{sO} \cdot \Delta i_{ds}^r(t) - R_{s,REF} \alpha_{cu} \Delta T_s^\circ(t) \cdot i_{dsO}^r + L_q \cdot i_{qsO}^r \cdot P_p \cdot \Delta \omega_m(t) \right. \\
        &\quad \left. + L_q \cdot P_p \cdot \omega_{mO} \cdot \Delta i_{qs}^r(t) \right] \\[10pt]
        %
        \Delta \dot{i}_{0s}^r(t) &= \frac{1}{L_{ls}} \left[ \Delta v_{0s}^r(t) - R_{sO} \cdot \Delta i_{0s}^r(t) - R_{s,REF} \alpha_{cu} \Delta T_s^\circ(t) \cdot i_{0s}^r \right] \\[10pt]
        %
        \Delta \dot{T}_s(t) &= \frac{1}{C_{ts}} \left[ \frac{3}{2} \cdot R_{sO} \cdot (2 \cdot i_{qsO}^r \cdot \Delta i_{qs}^r(t) + 2 \cdot i_{dsO}^r \cdot \Delta i_{ds}^r(t) + 4 \cdot i_{0sO}^r \cdot \Delta i_{0s}^r(t)) \right. \\
        &\quad + \frac{3 \alpha_{cu} R_{sREF}}{2} (i_{qsO}^r(t)^2 + i_{dsO}^r(t)^2 + 2i_{0sO}^r(t)^2) \cdot \Delta T_s(t) \\
        &\quad \left. - \frac{1}{R_{ts-amb}} (\Delta T_s(t) - \Delta T_{amb}(t)) \right]
    \end{aligned}
\right.
\end{equation*}
Donde \(R_{sO}=R_s(T^\circ_{sO})\)
\begin{equation*}
A_0 = \begin{bmatrix}
0 & 1 & 0 & 0 & 0 & 0 \\
-\frac{g k_l \cos\left(\frac{\theta_{mO}}{r}\right)}{J_{eq} r^2} & -\frac{b_{eq}}{J_{eq}} & \frac{3 P_p (\lambda_m + i_{dsO} (L_d - L_q))}{2 J_{eq}} & \frac{3 P_p i_{qsO} (L_d - L_q)}{2 J_{eq}} & 0 & 0 \\
0 & -\frac{P_p (\lambda_m + L_d i_{dsO})}{L_q} & -\frac{R_s}{L_q} & -\frac{L_d P_p \omega_{mO}}{L_q} & 0 & -\frac{R_{s,ref} \alpha_{cu} i_{qsO}}{L_q} \\
0 & \frac{L_q P_p i_{qsO}}{L_d} & \frac{L_q P_p \omega_{mO}}{L_d} & -\frac{R_s}{L_d} & 0 & -\frac{R_{s,ref} \alpha_{cu} i_{dsO}}{L_d} \\
0 & 0 & 0 & 0 & -\frac{R_s}{L_{ls}} & -\frac{R_{s,ref} \alpha_{cu} i_{0sO}}{L_{ls}} \\
0 & 0 & \frac{3 R_s i_{qsO}}{C_{ts}} & \frac{3 R_s i_{dsO}}{C_{ts}} & \frac{6 R_s i_{0sO}}{C_{ts}} & \Gamma_{O}
\end{bmatrix}
\end{equation*}

Donde la constante de disipación térmica acoplada $\Gamma_{th}$ se define como:
\begin{equation*}
\Gamma_{th} = -\frac{\frac{1}{R_{ts,amb}} - \frac{3 R_{s,ref} \alpha_{cu} (2 i_{0sO}^2 + i_{dsO}^2 + i_{qsO}^2)}{2}}{C_{ts}}
\end{equation*}
Las matrices de entrada manipulada y de perturbación quedan respectivamente:
\begin{equation*}
B_o^c = \begin{bmatrix}
0 & 0 & 0 \\
0 & 0 & 0 \\
\frac{1}{L_q} & 0 & 0 \\
0 & \frac{1}{L_d} & 0 \\
0 & 0 & \frac{1}{L_{ls}} \\
0 & 0 & 0
\end{bmatrix} ,
\end{equation*}

\begin{equation*}
B_o^d = \begin{bmatrix}
0 & 0 \\
-\frac{1}{J_{eq} r} & 0 \\
0 & 0 \\
0 & 0 \\
0 & 0 \\
0 & \frac{1}{C_{ts} R_{ts,amb}}
\end{bmatrix} 
\end{equation*}
El modelo global LPV toma la forma:
\begin{equation}
\begin{bmatrix}
\Delta \dot{\theta}_m(t) \\
\Delta \dot{\omega}_m(t) \\
\Delta \dot{i}_{qs}^r(t) \\
\Delta \dot{i}_{ds}^r(t) \\
\Delta \dot{i}_{0s}(t) \\
\Delta \dot{T}_s(t)
\end{bmatrix} = A_o \cdot 
\begin{bmatrix}
\Delta \theta_m(t) \\
\Delta \omega_m(t) \\
\Delta i_{qs}^r(t) \\
\Delta i_{ds}^r(t) \\
\Delta i_{0s}(t) \\
\Delta T_s(t)
\end{bmatrix} + B_o^c \cdot 
\begin{bmatrix}
\Delta v_{qs}^r(t) \\
\Delta v_{ds}^r(t) \\
\Delta v_{0s}(t)
\end{bmatrix} + B_o^d \cdot 
\begin{bmatrix}
\Delta T_{ld}(t) \\
\Delta T_{amb}(t)
\end{bmatrix} .
\end{equation}


\subsection{Linealización por Retroalimentación NL: Modelo LTI equivalente}
El objetivo de esta sección es obtener un modelo Lineal Invariante en el Tiempo (LTI) a partir de la planta no lineal, aplicando estrategias de Control Vectorial con campo orientado. Para lograr esto, se imponen restricciones matemáticas en las entradas manipuladas mediante retroalimentación de estado, buscando desacoplar los canales de flujo magnético y torque. Para el modelo LTI, se asume la dinámica térmica como lineal ($R_s = R_{s0}$ constante). Además, la no linealidad gravitacional de la carga se separa explícitamente de la matriz de estado, tratándola como una perturbación externa a ser compensada.

\subsubsection{Restricción o Ley de Control Mínima en el eje $d$}
Para lograr el desacoplamiento de flujo y torque, se impone el requisito de mantener la corriente del eje directo nula: $i_{ds}^r(t) \equiv 0$. Analizando la ecuación diferencial de dicho estado:
\begin{equation}
    \frac{di_{ds}^r(t)}{dt} = \frac{1}{L_d} \left[ v_{ds}^r(t) - R_{s0} i_{ds}^r(t) + L_q i_{qs}^r(t) \omega_r(t) \right]
\end{equation}

La \textbf{Ley de Control Mínima} requerida sobre la variable manipulada virtual $v_{ds}^r(t)$ para cancelar la fuerza contraelectromotriz cruzada y forzar la dinámica a cero (asumiendo la hipótesis de que el estado inicial es nulo, $i_{ds}^r(0) = 0$) es:
\begin{equation} \label{eq:vds_minima}
    v_{ds}^r(t) = -L_q \cdot i_{qs}^r(t) \cdot \omega_r(t)
\end{equation}

Al aplicar este voltaje, la dinámica del eje directo se anula y la ecuación de torque electromagnético se vuelve lineal y proporcional a $i_{qs}^r$.

\subsubsection{Ecuaciones Escalares y Matriciales del Modelo LTI (Nominal)}
Aplicando las leyes de control obtenidas \eqref{eq:vds_minima} y \eqref{eq:vqs_comp}, y asumiendo $i_{ds}^r(t) \equiv 0$, el sistema electromecánico global se reduce a un modelo LTI equivalente. El sistema de ecuaciones diferenciales resultante es:
\begin{equation} \label{eq:sis_escalar_LTI}
\begin{cases}
    \vspace{0.2cm}
    \dfrac{d\theta_m(t)}{dt} = \omega_m(t) \\
    \vspace{0.2cm}
    \dfrac{d\omega_m(t)}{dt} = \dfrac{1}{J_{eq}}\Big[\frac{3 P_p \lambda'_m}{2} i_{qs}(t) - b_{eq} \omega_m(t) - \frac{1}{r} T_{ld}(t) - \frac{g\cdot k_l}{r} \sin\left(\frac{\theta_m(t)}{r}\right)\Big] \\
    \dfrac{di_{qs}(t)}{dt} = \dfrac{1}{L_q}\Big[-P_p \lambda'_m \omega_m(t) - R_{sREF}\cdot i_{qs}(t) + v_{qs}(t)\Big]
\end{cases}
\end{equation}
Aunque se desprecie el efecto del subsistema térmico en la dinámica del modelo LTI, es interesante seguir estudiando la temperatura del estator en la simulación, por lo que también se considera la siguiente ecuación
\[\dfrac{dT_s(t)}{dt} = \dfrac{1}{C_{ts}}\Big[\frac{3}{2}R_{sREF}\cdot i_{qs}^2(t)-\frac{1}{R_{ts-amb}}\big(T_s(t) - T_{amb}(t)\big)\Big]\]
Definiendo el vector de estados nominal $\mathbf{x}_{LTI} = [\theta_m, \omega_m, i_{qs}]^T$, la representación matricial de estado (separando explícitamente las perturbaciones de carga y gravedad) resulta:

\begingroup
\renewcommand{\arraystretch}{1.8} 
\begin{equation}
\begin{bmatrix}
\dot{\theta}_m (t)\\
\dot{\omega}_m(t)\\
\dot{i}_{qs}(t)
\end{bmatrix}
= 
\underbrace{
\begin{bmatrix}
0 & 1 & 0 \\
0  & -\frac{b_{eq}}{J_{eq}} & \frac{3 P_p \lambda'_m}{2 J_{eq}}     \\
 0 &  -\frac{P_p \lambda'_m}{L_q} & -\frac{R_{s0}}{L_q}
\end{bmatrix}
}_{\mathbf{A}}
\begin{bmatrix}
\theta_m(t) \\
\omega_m (t)\\
i_{qs}(t) 
\end{bmatrix}
+
\underbrace{
\begin{bmatrix}
0  \\
0  \\
\frac{1}{L_q} 
\end{bmatrix}
}_{\mathbf{B_c}}
v_{qs}(t) 
+
\underbrace{
\begin{bmatrix}
0 \\
-\frac{1}{rJ_{eq}} \\
0
\end{bmatrix}
}_{\mathbf{B_d^1}}
T_{ld}(t)
+
\underbrace{
\begin{bmatrix}
    0 \\
    -\frac{g\cdot k_l}{rJ_{eq}} \\
    0
\end{bmatrix}
}_{\mathbf{B_d^2}}
\sin{\left(\frac{\theta_m(t)}{r}\right)}
\end{equation}\normalsize
\endgroup

Y la ecuación de salida midiendo la posición es:
\begin{equation}
\mathbf{y}(t) = \theta_m(t) = 
\underbrace{
\begin{bmatrix} 1 & 0 & 0 \end{bmatrix}
}_{\mathbf{C}_{LTI}} \mathbf{x}_{LTI} 
+ 
\underbrace{
\begin{bmatrix} 0 & 0 \end{bmatrix}
}_{\mathbf{D}_{LTI}} \mathbf{u}_{LTI}
\end{equation}



% Aquí insertas la imagen de tu diagrama de bloques del sistema LTI.
\begin{figure}[H]
    \centering
    \includegraphics[width=\textwidth]{imagenes/LTI equivalente_nominal.png}
    \caption{Diagrama de bloques de estado del modelo LTI equivalente nominal.}
    \label{fig:diagrama_LTI}
\end{figure}


\subsubsection{Dinámica residual equivalente y Modelo LTI Aumentado}
El modelo nominal desarrollado asume $i_{ds}^r(0) = 0$. En el caso general donde no se cumple esta hipótesis ($i_{ds}^r(0) \neq 0$), la inyección de la ley de control mínima \eqref{eq:vds_minima} en la ecuación del eje directo reduce la misma a una EDO autónoma de primer orden:
\begin{equation}
    \frac{di_{ds}^r(t)}{dt} = -\frac{R_{s0}}{L_d} i_{ds}^r(t) \implies i_{ds}^r(t) = i_{ds}^r(0) e^{-\frac{R_{s0}}{L_d} t}
\end{equation}

Esto demuestra una \textbf{dinámica residual} donde la corriente decae exponencialmente hacia cero. Como el acoplamiento cruzado en el eje $q$ ya fue cancelado matemáticamente por la ley \eqref{eq:vqs_comp}, esta dinámica transitoria ocurre de forma totalmente independiente del resto de los estados. 
Aun asi, esta volveria a acoplar temporalmente la corriente en eje directo con la velocidad, ya que:
\begin{equation*}
    v_{qs}^r(t) = L_q \dfrac{di_{qs}^r(t)}{dt} + R_s(t)i_{qs}^r(t)+\lambda_m'\omega_r(t)+  \underbrace{P_pL_di_{ds}^r(t)\omega_m(t)}_{\text{Acoplamiento}}
\end{equation*}
Para representar el sistema completo, se incorpora esta EDO al modelo LTI, definiendo el vector de estado aumentado $\mathbf{x}_{aug} = [\theta_m, \omega_m, i_{qs}, i_{ds}^r]^T$. 

\subsubsection{Ley de Control Complementaria en el eje $q$}
Si consideramos la dinámica residual analizada en el punto anterior, se puede ver que se producirá acoplamiento cruzado no lineal temporal en el eje de cuadratura debido al término $\omega_r L_d i_{ds}^r$. 

Para eliminar completamente esta dependencia y hacer que el eje $q$ sea lineal, se inyecta una compensación dinámica o \textbf{Ley de control complementaria}. La tensión virtual comandada será:
\begin{equation} \label{eq:vqs_comp}
    v_{qs}^{r*}(t) = v_{qs}(t) + \underbrace{\omega_r(t) L_d i_{ds}^r(t)}_{\text{Cancelación cruzada}}
\end{equation}
Donde $v_{qs}(t)$ pasa a ser la nueva entrada de control lineal del sistema.

El \textbf{Modelo LTI Aumentado} de cuarto orden resultante es:
\begin{equation} \label{eq:sis_escalar_LTI_aug}
\begin{cases}
    \vspace{0.2cm}
    \dfrac{d\theta_m(t)}{dt} = \omega_m(t) \\
    \vspace{0.2cm}
    \dfrac{d\omega_m(t)}{dt} = \dfrac{1}{J_{eq}}\Big[\frac{3 P_p \lambda'_m}{2} i_{qs}(t) - b_{eq} \omega_m(t) - \frac{1}{r} T_{ld}(t) - \frac{g\cdot k_l}{r} \sin\left(\frac{\theta_m(t)}{r}\right)\Big] \\
    \dfrac{di_{qs}(t)}{dt} = \dfrac{1}{L_q}\Big[-P_p \lambda'_m \omega_m(t) - R_{sREF}\cdot i_{qs}(t) + v_{qs}(t)\Big] \\
    \dfrac{di_{ds}(t)}{dt} = \dfrac{-R_{sREF}}{L_d}i_{ds}(t)\\
    \dfrac{dT_s(t)}{dt} = \dfrac{1}{C_{ts}}\Big[\frac{3}{2}R_{sREF}\cdot \big(i_{qs}^2(t)+i_{ds}^2(t)\big)-\frac{1}{R_{ts-amb}}\big(T_s(t) - T_{amb}(t)\big)\Big]
\end{cases}
\end{equation}
\begingroup
\renewcommand{\arraystretch}{1.8} 
\footnotesize\begin{equation}
\begin{bmatrix}
\dot{\theta}_m (t)\\
\dot{\omega}_m(t)\\
\dot{i}_{qs}(t)\\
\dot{i}_{ds}^r(t)
\end{bmatrix}
= 
\begin{bmatrix}
0 & 1 & 0 & 0 \\
0  & -\frac{b_{eq}}{J_{eq}} & \frac{3 P_p \lambda'_m}{2 J_{eq}} & 0 \\
 0 &  -\frac{P_p \lambda'_m}{L_q} & -\frac{R_{s0}}{L_q} & 0 \\
 0 & 0 & 0 & -\frac{R_{s0}}{L_d}
\end{bmatrix}
\begin{bmatrix}
\theta_m(t) \\
\omega_m (t)\\
i_{qs}(t) \\
i_{ds}^r(t)
\end{bmatrix}
+
\begin{bmatrix}
0  \\
0  \\
\frac{1}{L_q} \\
0
\end{bmatrix}
v_{qs}(t) 
+
\begin{bmatrix}
0 \\
-\frac{1}{rJ_{eq}} \\
0 \\
0
\end{bmatrix}
T_{ld}(t)
+
\begin{bmatrix}
    0 \\
    -\frac{g\cdot k_l}{rJ_{eq}} \\
    0 \\
    0
\end{bmatrix}
\sin{\left(\frac{\theta_m(t)}{r}\right)}
\end{equation}\normalsize
\endgroup
\subsubsection{Diagrama de bloques de estado (forma desagregada)}
A partir de las ecuaciones matriciales, se construye el diagrama de bloques del sistema linealizado equivalente.
\begin{figure}[H]
    \centering
    \includegraphics[width=\textwidth]{imagenes/LTI equivalente_aumentado.png}
    \caption{Diagrama de bloques de estado del modelo LTI equivalente aumentado.}
    \label{fig:diagrama_LTI_aumentado}
\end{figure}
\subsubsection{Implementación en el modelo global NL}
La implementación física de esta Linealización por Retroalimentación requiere un \textit{controlador parcial} acoplado a la planta NL original. 
\begin{enumerate}
    \item Recibe la acción de control lineal generada por un controlador de movimiento ($v_{qs}$).
    \item Calcula los voltajes de desacoplamiento interno en tiempo real utilizando las lecturas de los sensores idealizados ($\mathbf{i}_{abcs} \rightarrow \mathbf{i}_{qd0s}^r$ y $\omega_m$), aplicando las Ecuaciones \eqref{eq:vds_minima} y \eqref{eq:vqs_comp}.
    \item Estas señales resultantes ($v_{qs}^{r*}$ y $v_{ds}^{r*}$) pasan por un bloque de Transformación de Park Inversa ($\theta_r = P_p \theta_m$).
    \item Finalmente, se envían los voltajes calculados $\mathbf{v}_{abcs}^*$ al inversor, el cual excita físicamente la planta NL.
\end{enumerate}
\begin{figure}[H]
    \centering
    \includegraphics[width=\textwidth]{imagenes/NL con realimentacion NL.png}
    \caption{Diagrama de bloques de implementación de la realimentación sobre la planta no lineal.}
    \label{fig:diagrama_realimentacion_nl}
\end{figure}
\subsection{Comparación del modelo LPV con el modelo LTI aumentado}
\subsubsection{Comparación: LTI Aumentado vs. LPV con \(I_{ds0}^r \equiv 0\)}

Si bien ambos modelos buscan representar la dinámica del sistema mediante ecuaciones de estado lineales, difieren en su desarrollo y en cómo manejan las restricciones del sistema físico. Las principales diferencias y similitudes se detallan a continuación:

\begin{enumerate}
    \item \textbf{Naturaleza de la linealización:}
    \begin{itemize}
        \item \textit{Modelo LPV:} Se obtiene mediante una expansión en Series de Taylor truncada a primer orden. Es una aproximación de ``pequeña señal'' \cite{franklin_feedback}, válida únicamente en un entorno cercano al punto de equilibrio $(x_0, u_0)$.
        \item \textit{Modelo LTI Aumentado:} Se obtiene mediante Linealización por Retroalimentación NL. Es un modelo algebraicamente exacto sin aproximaciones de Taylor, logrado inyectando tensiones virtuales calculadas dinámicamente para cancelar las no linealidades de la planta.
    \end{itemize}
    
    \item \textbf{Tratamiento de la carga gravitacional (No linealidad mecánica):}
    \begin{itemize}
        \item \textit{Modelo LPV:} El torque gravitacional se linealiza, introduciendo un término de ``rigidez dinámica'' en la matriz de estado del sistema ($A_{2,1} \propto \cos(\theta_0/r)$). Esto provoca que los polos mecánicos cambien drásticamente según la posición geométrica del brazo.
        \item \textit{Modelo LTI Aumentado:} El término gravitacional completo $\sin(\theta_m/r)$ se extrae de la matriz $\mathbf{A}$ y se modela como una perturbación exógena que ingresa al sistema a través de la matriz $\mathbf{B_d^2}$. En consecuencia, la matriz de estado del LTI nominal ($A_{2,1} = 0$) permanece invariante frente a la posición.
    \end{itemize}

    \item \textbf{Desacoplamiento del Subsistema Térmico:}
    \begin{itemize}
        \item \textit{Modelo LPV:} Conserva la influencia de la temperatura como un parámetro de operación. La resistencia $R_{s0}$ dependerá del punto de equilibrio térmico $T_{s0}$ seleccionado para el cálculo del Jacobiano.
        \item \textit{Modelo LTI Aumentado:} Asume un desacoplamiento térmico completo, considerando a $R_s$ como una constante rígida, independizando al subsistema electromecánico de la evolución de la temperatura durante los transitorios rápidos.
    \end{itemize}

    \item \textbf{Similitud estructural al forzar $I_{ds0}^r \equiv 0$:} \\
    A pesar de las diferencias conceptuales, si evaluamos la matriz Jacobiana del modelo LPV forzando el punto de operación $I_{ds0}^r = 0$, observamos una gran convergencia hacia la estructura del modelo LTI Aumentado. 
    
    En el LPV original, los acoplamientos cruzados en la ecuación de tensión y torque generan términos en la matriz dependientes del producto cruzado de los estados. Específicamente, en la derivada de la fuerza contraelectromotriz respecto a la velocidad aparece el término $L_d I_{ds0}^r$. Si imponemos que $I_{ds0}^r = 0$:
    \begin{itemize}
        \item El acoplamiento cruzado de la velocidad sobre el eje $q$ se vuelve puramente proporcional al flujo de los imanes ($\lambda'_m$), coincidiendo exactamente con el término $A_{3,2}$ del LTI Aumentado.
        \item La contribución del torque de reluctancia desaparece en la linealización, quedando el torque dependiendo de forma puramente lineal de $I_{qs0}^r$, al igual que en el modelo LTI.
    \end{itemize}
    
    Por lo que, podemos decir que el Modelo LTI Aumentado puede interpretarse estructuralmente como el caso ideal de un modelo LPV forzado perpetuamente a operar sobre el punto $I_{ds0}^r = 0$, habiendo extraído las variaciones gravitacionales como perturbaciones de entrada.
\end{enumerate}
\subsubsection{Estudio del comportamiento del sistema sometido a variaciones de $i_{ds}^r$}
Al liberar la restricción $I_{ds0}^r = 0$, la corriente del eje directo pasa a ser un parámetro de sintonía que modifica tanto la capacidad estática de la máquina como sus propiedades dinámicas locales (autovalores del Jacobiano). El análisis de esta variación se divide en un enfoque físico-estático y uno dinámico.

\paragraph{1. Análisis Físico: Reforzamiento y Debilitamiento de Campo} \mbox{}\\
La inyección de una corriente continua en el eje directo altera el flujo magnético total de la máquina. El impacto directo se observa en la ecuación del par electromagnético en estado estacionario:
\begin{equation} \label{eq:torque_ids}
    T_m = \frac{3}{2} P_p \Big[ \lambda'_m + (L_d - L_q) I_{ds0}^r \Big] I_{qs0}^r
\end{equation}
Dependiendo del signo de la corriente inyectada, se distinguen dos regímenes de operación:
\begin{itemize}
    \item \textbf{Reforzamiento de campo ($I_{ds0}^r > 0$):} Para motores con polos salientes donde $L_d > L_q$, una corriente positiva genera un torque de reluctancia que se suma al torque de los imanes permanentes. Esto aumenta el torque total por amperio inyectado. Sin embargo, al reforzarse el flujo, la fuerza contraelectromotriz aumenta, lo que reduce la velocidad máxima alcanzable antes de saturar el inversor.\cite{krause_machinery}
    \item \textbf{Debilitamiento de campo ($I_{ds0}^r < 0$):} Una corriente negativa se opone al campo magnético principal de los imanes ($\lambda'_m$). Si bien esto disminuye la capacidad de producción de torque para una misma $I_{qs0}^r$, tiene el efecto beneficioso de reducir drásticamente la fuerza contraelectromotriz. Esto permite que el motor alcance velocidades de giro mucho mayores operando dentro del límite de tensión del inversor.\cite{krause_machinery}
\end{itemize}
\begin{figure}
    \centering
    \includegraphics[width=0.5\linewidth]{imagenes/polosaliente_cilindrico.png}
    \caption{Diagrama representativo de un rotor genérico de polos no-saliente (izq.) y un rotor de polos salientes (der.)\cite{estimation_excitation}}
    \label{fig:polos_saliente_cilindrico}
\end{figure}

\paragraph{2. Migración de Propiedades Dinámicas (Análisis de Autovalores)} \mbox{}\\
Para evaluar el impacto dinámico, se analiza la migración de los polos del sistema LPV (los autovalores de la matriz $\mathbf{A}_0$). Se fija un punto de operación nominal de velocidad ($\omega_{m0}$) y torque ($I_{qs0}^r$), y se realiza un barrido del parámetro $I_{ds0}^r$ desde la zona de debilitamiento hasta la zona de reforzamiento de campo. 

Sustituyendo estos valores en la matriz Jacobiana $\mathbf{A}_0$ analítica deducida previamente, se obtienen los autovalores para cada caso. Los resultados numéricos más relevantes se resumen en el Cuadro \ref{tab:migracion_polos}.


\begin{table}[H]
\centering
\caption{Debilitamiento de campo ($I_{ds0}^r<0$).}
\label{tab:migracion_polos}
\begin{tabular}{|c|c|c|c|c|c|}
\hline
$I_{ds0}^r$ [A] & $\lambda_1$ & $\lambda_2$ & $\lambda_3$ & $\lambda_4$ & $\lambda_5$ \\
\hline
-1.2 & -1275.00 & -141.30 & $-94.74 - j\,23.38$ & $-94.74 + j\,23.38$ & -0.2657 \\
-1.0 & -1275.00 & -145.62 & $-92.61 - j\,47.75$ & $-92.61 + j\,47.75$ & -0.2261 \\
-0.8 & -1275.00 & -147.64 & $-91.62 - j\,62.75$ & $-91.62 + j\,62.75$ & -0.1963 \\
-0.6 & -1275.00 & -148.88 & $-91.01 - j\,74.74$ & $-91.01 + j\,74.74$ & -0.1731 \\
-0.4 & -1275.00 & -149.74 & $-90.59 - j\,85.10$ & $-90.59 + j\,85.10$ & -0.1544 \\
\hline
\end{tabular}
\end{table}

\begin{table}[H]
\centering
\caption{Reforzamiento de campo ($I_{ds0}^r>0$).}
\label{tab:migracion_polos_reforzamiento}
\begin{tabular}{|c|c|c|c|c|c|}
\hline
$I_{ds0}^r$ [A] & $\lambda_1$ & $\lambda_2$ & $\lambda_3$ & $\lambda_4$ & $\lambda_5$ \\
\hline
0.2 & -1275.00 & -151.25 & $-89.85 - j\,111.05$ & $-89.85 + j\,111.05$ & -0.1157 \\
0.4 & -1275.00 & -151.57 & $-89.70 - j\,118.65$ & $-89.70 + j\,118.65$ & -0.1066 \\
0.6 & -1275.00 & -151.83 & $-89.57 - j\,125.90$ & $-89.57 + j\,125.90$ & -0.0986 \\
0.8 & -1275.00 & -152.05 & $-89.46 - j\,132.85$ & $-89.46 + j\,132.85$ & -0.0917 \\
1.0 & -1275.00 & -152.24 & $-89.37 - j\,139.56$ & $-89.37 + j\,139.56$ & -0.0857 \\
\hline
\end{tabular}
\end{table}

A partir de la evaluación numérica de la matriz Jacobiana, se desprenden las siguientes conclusiones sobre la migración de propiedades:
\begin{enumerate}
    \item \textbf{Frecuencia de oscilación:} La variación de $I_{ds0}^r$ afecta drásticamente a los términos cruzados electromecánicos en la matriz $\mathbf{A}_0$. Esto se traduce en una variación significativa de la parte imaginaria de los polos dominantes complejos conjugados. A medida que se pasa de debilitamiento a reforzamiento (aumenta $I_{ds0}^r$), la frecuencia natural de oscilación del sistema aumenta.
    \item \textbf{Estabilidad global:} La parte real de los polos dominantes complejos conjugados (que dicta el amortiguamiento absoluto) se mantiene prácticamente inalterada ante los cambios de $I_{ds0}^r$. Esto demuestra que inyectar corriente en el eje directo no compromete severamente la estabilidad a lazo abierto del sistema.
    \item \textbf{Dinámica rápida:} Los polos puramente reales asociados a la dinámica eléctrica rápida y a la dinámica térmica lenta permanecen virtualmente invariantes frente a los cambios en el punto de operación de $I_{ds0}^r$.
\end{enumerate}

% NOTA PARA TI: Puedes incluir aquí la gráfica (mapa de polos y ceros) generada por MATLAB superponiendo los polos para distintos I_ds0 (como la Figura 14 de tu referencia).
\begin{figure}[H]
    \centering
    % \includegraphics[width=0.8\textwidth]{ruta/migracion_polos.png}
    \caption{Mapa de migración de polos del modelo LPV variando $I_{ds0}^r$.}
    \label{fig:migracion_polos}
\end{figure}
\subsection{Funciones de Transferencia del Modelo LTI Aumentado}
A partir del modelo LTI equivalente aumentado de cuarto orden, se obtienen las funciones de transferencia desde ambas entradas ---la acción de control $v_{qs}(t)$ y la perturbación de carga $T_{ld}(t)$--- hacia la salida de posición $\theta_m(t)$.

\subsubsection{Estado inicial considerado}
El modelo LTI aumentado es un modelo \textbf{incremental} alrededor de un punto de equilibrio. El estado inicial considerado es:
\begin{equation}
    \mathbf{x}_{aug}(0) = \begin{bmatrix} \theta_m(0) \\ \omega_m(0) \\ i_{qs}(0) \\ i_{ds}^r(0) \end{bmatrix} = \begin{bmatrix} 0 \\ 0 \\ 0 \\ 0 \end{bmatrix}
\end{equation}
Además, dado que la gravedad fue extraída de la matriz de estado y se asume compensada mediante \textit{feedforward}, el término $\mathbf{B_d^2}\sin(\theta_m/r)$ no participa en la obtención de las funciones de transferencia lineales. La resistencia del estator se considera constante ($R_s = R_{s0}$) por el desacoplamiento térmico.

\subsubsection{Obtención de las funciones de transferencia}
Las funciones de transferencia se calculan mediante la expresión general:
\begin{equation}
    G(s) = \mathbf{C}(s\mathbf{I} - \mathbf{A})^{-1}\mathbf{B}
\end{equation}

Dado que el estado $i_{ds}^r$ está completamente desacoplado (como se demostrará en la sección siguiente), es suficiente trabajar con el subsistema nominal de tercer orden. Aplicando la transformada de Laplace (con condiciones iniciales nulas) a las ecuaciones escalares del modelo nominal:
\begin{equation}
\begin{cases}
    s\,\Theta_m(s) = \Omega_m(s) \\[6pt]
    s\,\Omega_m(s) = \dfrac{1}{J_{eq}}\left[\dfrac{3 P_p \lambda'_m}{2} I_{qs}(s) - b_{eq}\,\Omega_m(s) - \dfrac{1}{r}\,T_{ld}(s)\right] \\[6pt]
    s\,I_{qs}(s) = \dfrac{1}{L_q}\left[-P_p \lambda'_m\,\Omega_m(s) - R_{s0}\,I_{qs}(s) + V_{qs}(s)\right]
\end{cases}
\end{equation}

De la tercera ecuación se despeja $I_{qs}(s)$:
\begin{equation}\label{eq:iqs_despejada}
    I_{qs}(s) = \frac{V_{qs}(s) - P_p \lambda'_m\,\Omega_m(s)}{L_q\,s + R_{s0}}
\end{equation}

Sustituyendo \eqref{eq:iqs_despejada} en la segunda ecuación y agrupando:
\begin{equation}
    \left(J_{eq}\,s + b_{eq}\right)\Omega_m(s) = \frac{3 P_p \lambda'_m}{2} \cdot \frac{V_{qs}(s) - P_p \lambda'_m\,\Omega_m(s)}{L_q\,s + R_{s0}} - \frac{1}{r}\,T_{ld}(s)
\end{equation}

Multiplicando ambos miembros por $(L_q\,s + R_{s0})$ y reordenando:
\begin{equation}
    \left[(J_{eq}\,s + b_{eq})(L_q\,s + R_{s0}) + \frac{3 P_p^2 \lambda_m'^2}{2}\right]\Omega_m(s) = \frac{3 P_p \lambda'_m}{2}\,V_{qs}(s) - \frac{L_q\,s + R_{s0}}{r}\,T_{ld}(s)
\end{equation}

Finalmente, dado que $\Theta_m(s) = \Omega_m(s)/s$, y expandiendo el término entre corchetes se obtiene el denominador común:
\begin{equation}
    \Delta(s) = s\left[J_{eq} L_q\,s^2 + (b_{eq} L_q + J_{eq} R_{s0})\,s + b_{eq} R_{s0} + \frac{3 P_p^2 \lambda_m'^2}{2}\right]
\end{equation}

donde el factor $s$ proviene de la integración $\Theta_m = \Omega_m / s$. Dividiendo numerador y denominador por $J_{eq} L_q$ se obtienen las expresiones finales.

\vspace{0.3cm}

El polinomio característico del modelo aumentado es:
\begin{equation}
    \det(s\mathbf{I} - \mathbf{A}_{aug}) = \underbrace{\left(s + \frac{R_{s0}}{L_d}\right)}_{\text{Polo de } i_{ds}^r} \cdot \underbrace{\left[ s^3 + \left(\frac{b_{eq}}{J_{eq}} + \frac{R_{s0}}{L_q}\right)s^2 + \left(\frac{b_{eq} R_{s0}}{J_{eq} L_q} + \frac{3 P_p^2 \lambda_m'^2}{2 J_{eq} L_q}\right)s \right]}_{\Delta(s) \text{ --- Polinomio del subsistema nominal}}
\end{equation}

Sin embargo, al evalar la resolvente $(s\mathbf{I}-\mathbf{A})^{-1}$ y proyectar sobre las columnas de entrada $\mathbf{B_c}$ y $\mathbf{B_d}$, el factor $(s + R_{s0}/L_d)$ se \textbf{cancela} en ambas funciones de transferencia. Esto se debe a que el estado $i_{ds}^r$ no es excitado por ninguna de las dos entradas ($B_c(4) = B_d(4) = 0$) y no es observado por la salida ($C(4) = 0$). Por tanto, las funciones de transferencia resultan de \textbf{orden 3}, con denominador $\Delta(s)$:

\vspace{0.3cm}
\textbf{Función de transferencia de control} ($v_{qs} \to \theta_m$):
\begin{equation} \label{eq:tf_vqs}
\boxed{
    G_1(s) = \frac{\theta_m(s)}{v_{qs}(s)} = \frac{\dfrac{3 P_p \lambda'_m}{2 J_{eq} L_q}}{s^3 + \left(\dfrac{b_{eq}}{J_{eq}} + \dfrac{R_{s0}}{L_q}\right)s^2 + \left(\dfrac{b_{eq} R_{s0}}{J_{eq} L_q} + \dfrac{3 P_p^2 \lambda_m'^2}{2 J_{eq} L_q}\right)s}
}
\end{equation}

\vspace{0.3cm}
\textbf{Función de transferencia de perturbación} ($T_{ld} \to \theta_m$):
\begin{equation} \label{eq:tf_Tld}
\boxed{
    G_2(s) = \frac{\theta_m(s)}{T_{ld}(s)} = \frac{-\dfrac{1}{r J_{eq}}\left(s + \dfrac{R_{s0}}{L_q}\right)}{s^3 + \left(\dfrac{b_{eq}}{J_{eq}} + \dfrac{R_{s0}}{L_q}\right)s^2 + \left(\dfrac{b_{eq} R_{s0}}{J_{eq} L_q} + \dfrac{3 P_p^2 \lambda_m'^2}{2 J_{eq} L_q}\right)s}
}
\end{equation}

Ambas funciones de transferencia comparten el mismo polinomio característico $\Delta(s)$, el cual presenta un polo en el origen ($s=0$) correspondiente a la integración inherente de velocidad a posición. Esto indica que el sistema a lazo abierto es de \textbf{tipo 1}: sigue rampas de posición con error estacionario finito, pero no alcanza error nulo ante entrada escalón sin realimentación.

Notar que $G_2(s)$ posee un \textbf{cero} en $s = -R_{s0}/L_q$, que coincide con el polo eléctrico del eje $q$. Esto se interpreta físicamente como: ante una perturbación de torque de carga, la corriente $i_{qs}$ reacciona inicialmente absorbiendo parte del impacto, atenuando la transmisión del disturbio hacia la posición en frecuencias altas.

\subsubsection{Estado interno no reflejado en las funciones de transferencia}
El estado $i_{ds}^r$ (cuarto estado del modelo aumentado) \textbf{no aporta ni se ve reflejado} en ninguna de las funciones de transferencia obtenidas. Formalmente, esto se verifica analizando las propiedades estructurales del sistema:

\begin{enumerate}
    \item \textbf{No controlabilidad:} Ninguna de las entradas ($v_{qs}$, $T_{ld}$) excita al estado $i_{ds}^r$, ya que las componentes correspondientes de las matrices de entrada son nulas: $B_c(4) = 0$ y $B_d(4) = 0$.
    \item \textbf{No observabilidad:} La salida $\theta_m$ no depende de $i_{ds}^r$, ya que $C(4) = 0$.
    \item \textbf{Desacoplamiento dinámico:} La cuarta fila de $\mathbf{A}$ solo contiene el término $-R_{s0}/L_d$ en la diagonal, sin acoplamientos con los demás estados. La dinámica de $i_{ds}^r$ es completamente autónoma.
\end{enumerate}

Consecuentemente, el polo $\lambda_4 = -R_{s0}/L_d$ del sistema en espacio de estados se cancela algebraicamente al convertir a función de transferencia. La dinámica de $i_{ds}^r$ es puramente una \textbf{respuesta libre}:
\begin{equation}
    i_{ds}^r(t) = i_{ds}^r(0) \, e^{-\frac{R_{s0}}{L_d} t}
\end{equation}
que solo se manifiesta si la condición inicial es no nula, decayendo exponencialmente con constante de tiempo $\tau_d = L_d / R_{s0}$. Esto confirma que el \textbf{modelo nominal de tercer orden} $\mathbf{x}_{LTI} = [\theta_m, \omega_m, i_{qs}]^T$ es suficiente para caracterizar completamente la relación entrada-salida del sistema.

\subsection{Análisis de Estabilidad a Lazo Abierto}

\subsubsection{Autovalores del modelo aumentado}
Los autovalores del sistema se obtienen a partir del polinomio característico de la matriz $\mathbf{A}_{aug}$ (4$\times$4). Como se demostró en la sección anterior, este se factoriza en dos términos independientes:
\begin{equation}
    \det(s\mathbf{I} - \mathbf{A}_{aug}) = \underbrace{\left(s + \frac{R_{s0}}{L_d}\right)}_{\lambda_4} \cdot \underbrace{s}_{\lambda_1} \cdot \underbrace{\left(s^2 + a_2\,s + a_1\right)}_{\lambda_{2,3}}
\end{equation}
donde los coeficientes del polinomio cuadrático son:
\begin{equation}
    a_2 = \frac{b_{eq}}{J_{eq}} + \frac{R_{s0}}{L_q}, \qquad a_1 = \frac{b_{eq} R_{s0}}{J_{eq} L_q} + \frac{3 P_p^2 \lambda_m'^2}{2 J_{eq} L_q}
\end{equation}

Evaluando numéricamente con los parámetros nominales ($m_l = 0$, $J_{l,nom} = 0{,}0833$ kg$\cdot$m$^2$):

\begin{table}[H]
\centering
\caption{Autovalores del modelo LTI aumentado a lazo abierto (parámetros nominales).}
\label{tab:autovalores_LA}
\renewcommand{\arraystretch}{1.4}
\begin{tabular}{|c|c|c|c|c|}
\hline
\textbf{Autovalor} & \textbf{Valor} & \textbf{Subsistema} & $\boldsymbol{\omega_n}$ \textbf{[rad/s]} & $\boldsymbol{\zeta}$ \\
\hline
$\lambda_1$ & $0$ & Integrador (posición) & --- & --- \\
\hline
$\lambda_{2,3}$ & $-88{,}5 \pm j\,150{,}0$ & Electromecánico & $174{,}1$ & $0{,}508$ \\
\hline
$\lambda_4$ & $-154{,}5$ & Residual $i_{ds}^r$ & --- & $\tau_d = 6{,}5$ ms \\
\hline
\end{tabular}
\end{table}

\paragraph{Identificación física de cada autovalor:} \mbox{}

\begin{itemize}
    \item \textbf{$\lambda_1 = 0$ (Polo en el origen):} Corresponde a la relación cinemática $\dot{\theta}_m = \omega_m$. La posición es la integral de la velocidad, por lo que el sistema posee un \textbf{integrador puro}. Esto implica que ante una entrada constante, la posición crece indefinidamente sin realimentación.

    \item \textbf{$\lambda_{2,3} = -88{,}5 \pm j\,150{,}0$ (Par complejo conjugado):} Representan el \textbf{modo electromecánico acoplado} entre la corriente $i_{qs}$ y la velocidad $\omega_m$. La parte imaginaria indica un comportamiento oscilatorio subamortiguado ($\zeta \approx 0{,}51$), producto del intercambio de energía entre el campo magnético ($L_q$, $\lambda'_m$) y la inercia mecánica ($J_{eq}$). La frecuencia de oscilación amortiguada es $\omega_d = 150{,}0$ rad/s ($\approx 23{,}9$ Hz), y la envolvente decae con constante de tiempo $\tau_{em} = 1/88{,}5 \approx 11{,}3$ ms.

    \item \textbf{$\lambda_4 = -R_{s0}/L_d = -154{,}5$ rad/s (Polo real negativo):} Corresponde a la dinámica residual autónoma de $i_{ds}^r$. Como se demostró, este polo \textbf{no aparece} en las funciones de transferencia. Su constante de tiempo es $\tau_d = L_d/R_{s0} \approx 6{,}5$ ms.
\end{itemize}

\subsubsection{Ceros de las funciones de transferencia}

\begin{itemize}
    \item \textbf{$G_1(s) = \theta_m / v_{qs}$:} No posee ceros finitos. Es un sistema de \textbf{fase mínima} puro de tipo 1 con exceso polar de 3. La respuesta ante un escalón de tensión será una rampa modulada por la dinámica oscilatoria de los polos electromecánicos.

    \item \textbf{$G_2(s) = \theta_m / T_{ld}$:} Posee un cero real negativo en:
    \begin{equation}
        z_1 = -\frac{R_{s0}}{L_q} = -175{,}9 \text{ rad/s} \qquad (\tau_z = 5{,}7 \text{ ms})
    \end{equation}
    Este cero se ubica muy cerca del polo electromecánico dominante. Físicamente, representa que el subsistema eléctrico ($L_q$, $R_{s0}$) ``absorbe'' parcialmente las perturbaciones de torque antes de que estas se transmitan completamente a la posición, actuando como un filtro natural de alta frecuencia sobre la perturbación.
\end{itemize}

\subsubsection{Estabilidad}

\textbf{Estabilidad interna (Lyapunov):} \mbox{} \\
El modelo LTI aumentado posee un autovalor en $s = 0$ (simple) y los restantes con parte real estrictamente negativa. Por tanto, el sistema es \textbf{marginalmente estable} (o estable en el sentido de Lyapunov, pero \textbf{no} asintóticamente estable). Una perturbación en la posición no se atenúa: el sistema ``recuerda'' cualquier desplazamiento permanentemente.

\textbf{Estabilidad BIBO (entrada-salida acotada):} \mbox{} \\
Para la salida $\theta_m$:
\begin{itemize}
    \item Desde $v_{qs}$: \textbf{No es BIBO estable}. El polo en el origen de $G_1(s)$ implica que un escalón constante de tensión produce una posición que crece sin cota (rampa).
    \item Desde $T_{ld}$: \textbf{No es BIBO estable}. Análogamente, un torque de perturbación constante genera una deriva indefinida de la posición.
\end{itemize}

\textbf{Estabilidad parcial:} \mbox{} \\
Si se analiza únicamente el subsistema de velocidad y corriente ($\omega_m$, $i_{qs}$), los polos electromecánicos $\lambda_{2,3}$ tienen parte real negativa, garantizando que las \textbf{variables de velocidad y corriente convergen} a valores estacionarios finitos ante entradas acotadas. Solo la posición es la variable inestable (integradora).

Esta estructura es característica de los sistemas de \textbf{regulación de posición}: se requiere necesariamente un lazo cerrado de realimentación para estabilizar la posición. El controlador PID que se diseñará en la sección 5.2 cumplirá exactamente esta función.

\begin{figure}[H]
    \centering
    \includegraphics[width=\textwidth]{imagenes/polos_ceros_LA.png}
    \caption{Mapa de polos y ceros a lazo abierto para ambas funciones de transferencia del modelo LTI aumentado. Se indica con $\times$ roja el polo cancelado de $i_{ds}^r$.}
    \label{fig:polos_ceros_LA}
\end{figure}

\subsubsection{Migración de propiedades ante variación de \texorpdfstring{$R_s$}{Rs}}
La resistencia del estator varía con la temperatura según $R_s(T_s) = R_{s0}(1 + \alpha_{Cu}(T_s - T_{s,ref}))$. Para evaluar su efecto sobre la dinámica, se realiza un barrido en el rango $[0{,}8\,R_{s0},\; 1{,}2\,R_{s0}]$.

\begin{table}[H]
\centering
\caption{Migración de polos y ceros ante variación de $R_s$.}
\label{tab:migracion_Rs}
\renewcommand{\arraystretch}{1.2}
\footnotesize
\begin{tabular}{|c|c|c|c|c|}
\hline
$R_s$ [$\Omega$] & Polos $\lambda_{2,3}$ & Cero $G_2$ & $\omega_n$ [rad/s] & $\zeta$ \\
\hline
0.816 & $-70{,}9 \pm j\,158{,}9$ & $-140{,}69$ & $174{,}0$ & $0{,}407$ \\
0.850 & $-73{,}8 \pm j\,157{,}6$ & $-146{,}55$ & $174{,}0$ & $0{,}424$ \\
0.884 & $-76{,}8 \pm j\,156{,}2$ & $-152{,}41$ & $174{,}0$ & $0{,}441$ \\
0.918 & $-79{,}7 \pm j\,154{,}7$ & $-158{,}28$ & $174{,}0$ & $0{,}458$ \\
0.952 & $-82{,}6 \pm j\,153{,}2$ & $-164{,}14$ & $174{,}1$ & $0{,}475$ \\
0.986 & $-85{,}6 \pm j\,151{,}6$ & $-170{,}00$ & $174{,}1$ & $0{,}491$ \\
\textbf{1.020} & $\mathbf{-88{,}5 \pm j\,149{,}9}$ & $\mathbf{-175{,}86}$ & $\mathbf{174{,}1}$ & $\mathbf{0{,}508}$ \\
1.054 & $-91{,}4 \pm j\,148{,}2$ & $-181{,}72$ & $174{,}1$ & $0{,}525$ \\
1.088 & $-94{,}3 \pm j\,146{,}4$ & $-187{,}59$ & $174{,}1$ & $0{,}542$ \\
1.122 & $-97{,}3 \pm j\,144{,}5$ & $-193{,}45$ & $174{,}2$ & $0{,}559$ \\
1.156 & $-100{,}2 \pm j\,142{,}5$ & $-199{,}31$ & $174{,}2$ & $0{,}575$ \\
1.190 & $-103{,}1 \pm j\,140{,}4$ & $-205{,}17$ & $174{,}2$ & $0{,}592$ \\
1.224 & $-106{,}1 \pm j\,138{,}2$ & $-211{,}03$ & $174{,}2$ & $0{,}609$ \\
\hline
\end{tabular}
\end{table}

Del análisis se concluye:
\begin{itemize}
    \item La \textbf{frecuencia natural} $\omega_n$ permanece prácticamente constante ($\approx 174$ rad/s), ya que depende principalmente del acoplamiento electromecánico $P_p\lambda'_m$ y la inercia $J_{eq}$, no de $R_s$.
    \item El \textbf{amortiguamiento} $\zeta$ aumenta con $R_s$ (de $0{,}41$ a $0{,}61$), lo que se explica porque una mayor resistencia disipa más rápidamente la energía almacenada en el campo magnético.
    \item El \textbf{cero} de $G_2(s)$ migra proporcionalmente a $R_s$ ($z = -R_s/L_q$), desplazándose hacia la izquierda con el aumento de temperatura.
\end{itemize}

\begin{figure}[H]
    \centering
    \includegraphics[width=\textwidth]{imagenes/migracion_Rs.png}
    \caption{Migración de polos y ceros en el plano $s$ (izquierda) y variación de $\omega_n$ y $\zeta$ (derecha) ante cambios en $R_s$.}
    \label{fig:migracion_Rs}
\end{figure}

\subsubsection{Migración de propiedades ante variación de carga ($m_l$)}
La masa de carga útil $m_l \in [0,\; 1{,}5]$ kg modifica la inercia equivalente $J_{eq} = J_m + J_l/r^2$ y el coeficiente gravitacional $k_l$. Al aumentar $m_l$, la inercia crece y el acoplamiento electromecánico se ``diluye'', modificando los autovalores.

\begin{table}[H]
\centering
\caption{Migración de polos ante variación de masa de carga $m_l$.}
\label{tab:migracion_ml}
\renewcommand{\arraystretch}{1.2}
\footnotesize
\begin{tabular}{|c|c|c|c|c|}
\hline
$m_l$ [kg] & $J_{eq}$ [kg$\cdot$m$^2$] & Polos $\lambda_{2,3}$ & $\omega_n$ [rad/s] & $\zeta$ \\
\hline
\textbf{0.00} & $\mathbf{1{,}98 \times 10^{-5}}$ & $\mathbf{-88{,}5 \pm j\,149{,}9}$ & $\mathbf{174{,}1}$ & $\mathbf{0{,}508}$ \\
0.20 & $2{,}33 \times 10^{-5}$ & $-88{,}4 \pm j\,134{,}1$ & $160{,}6$ & $0{,}551$ \\
0.40 & $2{,}67 \times 10^{-5}$ & $-88{,}3 \pm j\,121{,}0$ & $149{,}8$ & $0{,}590$ \\
0.60 & $3{,}02 \times 10^{-5}$ & $-88{,}3 \pm j\,109{,}8$ & $140{,}9$ & $0{,}627$ \\
0.80 & $3{,}37 \times 10^{-5}$ & $-88{,}3 \pm j\,100{,}1$ & $133{,}5$ & $0{,}661$ \\
1.00 & $3{,}71 \times 10^{-5}$ & $-88{,}2 \pm j\,91{,}4$ & $127{,}1$ & $0{,}694$ \\
1.20 & $4{,}06 \times 10^{-5}$ & $-88{,}2 \pm j\,83{,}6$ & $121{,}5$ & $0{,}726$ \\
1.50 & $4{,}58 \times 10^{-5}$ & $-88{,}2 \pm j\,72{,}9$ & $114{,}4$ & $0{,}771$ \\
\hline
\end{tabular}
\end{table}

\begin{figure}[H]
    \centering
    \includegraphics[width=\textwidth]{imagenes/migracion_ml.png}
    \caption{Migración de polos ante variación de masa de carga $m_l$ (izquierda) y variación de $\omega_n$ y $\zeta$ (derecha).}
    \label{fig:migracion_ml}
\end{figure}

Del análisis se observa que al aumentar $m_l$:
\begin{itemize}
    \item La \textbf{frecuencia natural} $\omega_n$ disminuye, ya que la mayor inercia reduce la relación torque/inercia del sistema.
    \item El \textbf{amortiguamiento} $\zeta$ aumenta, porque la fricción equivalente $b_{eq}$ se mantiene mientras la frecuencia natural baja, incrementando la relación $\zeta = a_2/(2\omega_n)$.
    \item El polo en el origen y el cero de $G_2$ permanecen invariantes frente a $m_l$.
\end{itemize}

Estas variaciones justifican la necesidad de un controlador robusto capaz de mantener el desempeño dentro de todo el rango de incertidumbre de carga.

\subsubsection{Diagramas de Bode a lazo abierto}
\begin{figure}[H]
    \centering
    \includegraphics[width=\textwidth]{imagenes/bode_LA.png}
    \caption{Diagramas de Bode de las funciones de transferencia $G_1(s)$ y $G_2(s)$ a lazo abierto.}
    \label{fig:bode_LA}
\end{figure}

\subsection{Análisis de Observabilidad}
Se analiza si es posible reconstruir el vector de estado completo del modelo LTI aumentado a partir de la medición de una única salida escalar.

\subsubsection{Observabilidad desde \texorpdfstring{$\theta_m(t)$}{theta_m(t)} (sensor de posición)}
La matriz de observabilidad para el sistema $(\mathbf{A}_{aug}, \mathbf{C})$ con $\mathbf{C} = \begin{bmatrix} 1 & 0 & 0 & 0 \end{bmatrix}$ se construye como:
\begin{equation}
    \mathcal{O} = \begin{bmatrix} \mathbf{C} \\ \mathbf{C}\mathbf{A} \\ \mathbf{C}\mathbf{A}^2 \\ \mathbf{C}\mathbf{A}^3 \end{bmatrix}
\end{equation}

Evaluando fila por fila a partir de la estructura de $\mathbf{A}_{aug}$:
\begin{itemize}
    \item $\mathbf{C} = \begin{bmatrix} 1 & 0 & 0 & 0 \end{bmatrix}$ --- se observa $\theta_m$ directamente.
    \item $\mathbf{C}\mathbf{A} = \begin{bmatrix} 0 & 1 & 0 & 0 \end{bmatrix}$ --- se observa $\omega_m$ (derivada de $\theta_m$).
    \item $\mathbf{C}\mathbf{A}^2 = \begin{bmatrix} 0 & -\frac{b_{eq}}{J_{eq}} & \frac{3P_p\lambda'_m}{2J_{eq}} & 0 \end{bmatrix}$ --- combinación de $\omega_m$ e $i_{qs}$.
    \item $\mathbf{C}\mathbf{A}^3$: combinación que involucra $\omega_m$ e $i_{qs}$, pero \textbf{nunca} $i_{ds}^r$.
\end{itemize}

La cuarta columna de $\mathcal{O}$ es idénticamente nula, ya que $i_{ds}^r$ no acopla con ningún estado observable. Por lo tanto:
\begin{equation}
    \text{rango}(\mathcal{O}) = 3 < 4 = n
\end{equation}

\textbf{Conclusión:} El sistema \textbf{no es completamente observable} desde $\theta_m(t)$. El estado no observable es $i_{ds}^r$. Los tres estados restantes ($\theta_m$, $\omega_m$, $i_{qs}$) sí son observables, lo que permite diseñar un observador reducido para estimar $\omega_m$ e $i_{qs}$ a partir de la medición de posición.

Esto es consistente con el análisis de las funciones de transferencia: dado que $i_{ds}^r$ no se manifiesta en la salida $\theta_m$, es imposible reconstruirlo a partir de dicha medición.

\subsubsection{Observabilidad desde \texorpdfstring{$\omega_m(t)$}{omega_m(t)} (tacogenerador)}
Si en lugar de medir la posición se mide la velocidad angular, la matriz de salida es $\mathbf{C}_\omega = \begin{bmatrix} 0 & 1 & 0 & 0 \end{bmatrix}$. Evaluando:
\begin{itemize}
    \item $\mathbf{C}_\omega = \begin{bmatrix} 0 & 1 & 0 & 0 \end{bmatrix}$ --- se observa $\omega_m$.
    \item $\mathbf{C}_\omega\mathbf{A} = \begin{bmatrix} 0 & -\frac{b_{eq}}{J_{eq}} & \frac{3P_p\lambda'_m}{2J_{eq}} & 0 \end{bmatrix}$ --- combinación de $\omega_m$ e $i_{qs}$.
    \item $\mathbf{C}_\omega\mathbf{A}^2$: combinación de $\omega_m$ e $i_{qs}$.
    \item $\mathbf{C}_\omega\mathbf{A}^3$: combinación de $\omega_m$ e $i_{qs}$.
\end{itemize}

Nuevamente, la cuarta columna es nula ($i_{ds}^r$ sigue sin ser observable), pero además la primera columna también es nula: la posición $\theta_m$ nunca aparece en las potencias de $\mathbf{C}_\omega\mathbf{A}^k$ porque $\theta_m$ solo se acopla unidireccionalmente como integrador de $\omega_m$.

\begin{equation}
    \text{rango}(\mathcal{O}_\omega) = 2 < 4 = n
\end{equation}

\textbf{Conclusión:} Midiendo $\omega_m$ se pierden \textbf{dos} estados: $\theta_m$ e $i_{ds}^r$. La posición no es observable desde la velocidad porque la integración ``borra'' la condición inicial ---es decir, infinitas trayectorias de posición producen la misma velocidad. Por tanto, medir $\theta_m$ con un \textit{encoder} es estrictamente más informativo que medir $\omega_m$ con un tacogenerador para este sistema.

\subsection{Análisis de Controlabilidad}
Se analiza si es posible llevar el vector de estado completo a cualquier valor deseado mediante la entrada de control $v_{qs}(t)$, sin considerar la perturbación $T_{ld}(t)$.

\subsubsection{Controlabilidad desde \texorpdfstring{$v_{qs}(t)$}{v_qs(t)}}
La matriz de controlabilidad para el sistema $(\mathbf{A}_{aug}, \mathbf{B}_c)$ con $\mathbf{B}_c = \begin{bmatrix} 0 & 0 & 1/L_q & 0 \end{bmatrix}^T$ es:
\begin{equation}
    \mathcal{C} = \begin{bmatrix} \mathbf{B}_c & \mathbf{A}\mathbf{B}_c & \mathbf{A}^2\mathbf{B}_c & \mathbf{A}^3\mathbf{B}_c \end{bmatrix}
\end{equation}

Evaluando columna por columna:
\begin{itemize}
    \item $\mathbf{B}_c = \begin{bmatrix} 0 \\ 0 \\ 1/L_q \\ 0 \end{bmatrix}$ --- la entrada actúa directamente sobre $i_{qs}$.
    \item $\mathbf{A}\mathbf{B}_c = \begin{bmatrix} 0 \\ \frac{3P_p\lambda'_m}{2J_{eq}L_q} \\ -\frac{R_{s0}}{L_q^2} \\ 0 \end{bmatrix}$ --- la acción se propaga a $\omega_m$ a través del torque.
    \item $\mathbf{A}^2\mathbf{B}_c$: la acción alcanza $\theta_m$ (vía la cadena $i_{qs} \to \omega_m \to \theta_m$).
    \item $\mathbf{A}^3\mathbf{B}_c$: combinación de los tres estados del subsistema nominal.
\end{itemize}

La cuarta fila de todas las columnas de $\mathcal{C}$ es idénticamente nula, ya que la entrada $v_{qs}$ nunca llega al estado $i_{ds}^r$ ($B_c(4) = 0$) y ningún otro estado acopla hacia $i_{ds}^r$ (la fila 4 de $\mathbf{A}$ es diagonal). Por lo tanto:
\begin{equation}
    \text{rango}(\mathcal{C}) = 3 < 4 = n
\end{equation}

\textbf{Conclusión:} El sistema \textbf{no es completamente controlable} desde $v_{qs}(t)$. El estado no controlable es $i_{ds}^r$.

\subsubsection{Incorporación de una entrada adicional}
El estado $i_{ds}^r$ podría controlarse si se agrega una segunda entrada manipulable: la tensión del eje directo $v_{ds}^r(t)$. En el modelo aumentado, esto equivale a extender la matriz de entrada:
\begin{equation}
    \mathbf{B}_{ext} = \begin{bmatrix} \mathbf{B}_c & \mathbf{B}_{ds} \end{bmatrix} =
    \begin{bmatrix}
    0 & 0 \\ 0 & 0 \\ 1/L_q & 0 \\ 0 & 1/L_d
    \end{bmatrix}
\end{equation}

Con esta extensión, la matriz de controlabilidad alcanza rango 4 y el sistema pasa a ser completamente controlable. Sin embargo, en la práctica esta entrada $v_{ds}^r$ ya está siendo utilizada por la \textbf{ley de control mínima} \eqref{eq:vds_minima} para forzar $i_{ds}^r \to 0$. Es decir, la propia estrategia de control vectorial ya ``controla'' $i_{ds}^r$ llevándola exponencialmente a cero; simplemente no lo hace a través de la estructura de realimentación de estado convencional, sino mediante la inyección de una tensión de desacoplamiento.

\subsubsection{Resumen de propiedades estructurales}

\begin{table}[H]
\centering
\caption{Resumen de controlabilidad y observabilidad del modelo LTI aumentado.}
\label{tab:ctrl_obs}
\renewcommand{\arraystretch}{1.3}
\begin{tabular}{|l|c|c|c|c|}
\hline
\textbf{Estado} & \textbf{Controlable} ($v_{qs}$) & \textbf{Observable} ($\theta_m$) & \textbf{Observable} ($\omega_m$) & \textbf{En las FTs} \\
\hline
$\theta_m$   & Sí & Sí (directo) & No & Sí \\
$\omega_m$   & Sí & Sí           & Sí (directo) & Sí \\
$i_{qs}$     & Sí (directo) & Sí      & Sí & Sí \\
$i_{ds}^r$   & No & No            & No & No \\
\hline
\end{tabular}
\end{table}

El estado $i_{ds}^r$ es simultáneamente no controlable y no observable, confirmando su naturaleza de \textbf{modo oculto} del sistema. Este resultado es coherente con la cancelación polo-cero identificada en las funciones de transferencia y con el desacoplamiento estructural de la cuarta fila y columna de la matriz $\mathbf{A}_{aug}$.

%%%SECCION 516
\subsection{Simulación Dinámica: Modelo NL vs LTI Aumentado}
\label{sec:sim_dinamica}

Se realiza una simulación en el dominio del tiempo comparando la respuesta del \textbf{modelo no lineal completo} (6 estados, con leyes de control NL de desacoplamiento mínima y complementaria) contra el \textbf{modelo LTI aumentado} (4 estados), implementados en MATLAB/Simulink. La simulación se ejecuta para tres condiciones iniciales: $i_{ds}^r(0) \in \{0,\; +0{,}5,\; -0{,}5\}$ A, manteniendo el resto del estado inicial en cero y $T_{amb} = 20\,^\circ$C.

\subsubsection{Señales de excitación}

La consigna define una secuencia temporal de 10 eventos (escalones) que producen transitorios sucesivos. La combinación de ambas señales permite explorar los cuatro cuadrantes de operación.

\textbf{Pulso de tensión en eje $q$:}
\begin{equation}
v_{qs}^*(t) = \begin{cases}
0 \text{ V} & t < 0{,}1 \text{ s} \\
+19{,}596 \text{ V} & 0{,}1 \leq t < 0{,}7 \text{ s} \\
0 \text{ V} & 0{,}7 \leq t < 1{,}1 \text{ s} \\
-19{,}596 \text{ V} & 1{,}1 \leq t < 1{,}7 \text{ s} \\
0 \text{ V} & t \geq 1{,}7 \text{ s}
\end{cases}
\end{equation}

\textbf{Doble pulso de torque de carga:}
\begin{equation}
T_{ld}(t) = \begin{cases}
0 & t < 0{,}3 \text{ s} \\
+6{,}28 \text{ N.m} & 0{,}3 \leq t < 0{,}5 \text{ s} \\
-6{,}28 \text{ N.m} & 0{,}5 \leq t < 0{,}9 \text{ s} \\
0 & 0{,}9 \leq t < 1{,}3 \text{ s} \\
+6{,}28 \text{ N.m} & 1{,}3 \leq t < 1{,}5 \text{ s} \\
-6{,}28 \text{ N.m} & 1{,}5 \leq t < 1{,}9 \text{ s} \\
0 & t \geq 1{,}9 \text{ s}
\end{cases}
\end{equation}

\begin{figure}[H]
    \centering
    \includegraphics[width=\textwidth]{imagenes/sim_senales_entrada.png}
    \caption{Señales de excitación: consigna de tensión $v_{qs}^*(t)$ y torque de carga $T_{ld}(t)$.}
    \label{fig:sim_entradas}
\end{figure}

De la secuencia se identifican 10 transitorios consecutivos, separados por los instantes de conmutación $t_1 = 0{,}1$~s, $t_2 = 0{,}3$~s, $t_3 = 0{,}5$~s, $t_4 = 0{,}7$~s, $t_5 = 0{,}9$~s, $t_6 = 1{,}1$~s, $t_7 = 1{,}3$~s, $t_8 = 1{,}5$~s, $t_9 = 1{,}7$~s, $t_{10} = 1{,}9$~s. La primera mitad (Trans.~1--5) corresponde al pulso positivo de $v_{qs}^*$ y la segunda mitad (Trans.~6--10) a su simétrico negativo.

\subsubsection{Modelo en Simulink}

Se construyó un modelo en Simulink con dos subsistemas en paralelo: el modelo NL completo y el modelo LTI aumentado. Ambos reciben las mismas señales de excitación.

\begin{figure}[H]
    \centering
    \includegraphics[width=0.5\textwidth]{imagenes/sim_diagrama_bloques.png}
    \caption{Diagrama de bloques en Simulink: modelo NL desacoplado y LTI aumentado en paralelo frente al mismo Vq* consigna y Torque de Perturbacion.}
    \label{fig:simulink_diagrama}
\end{figure}

\subsubsection{Respuesta del estado interno (Sub-ítem a)}
\label{sec:sim_item_a}

Se presentan los resultados de la simulación para el caso nominal $i_{ds}^r(0) = 0$A, comparando las salidas del modelo NL (6 estados) y del modelo LTI aumentado (4 estados) ante las mismas señales de excitación. Los efectos de $i_{ds}^r(0) \neq 0$ se analizan por separado en el Sub-ítem~c. En las siguientes siulaciones se consideró $T_{amb}=20°$ y en el LTI aumentado se consideró R_s promedio entre 

\textbf{Posición angular $\theta_m(t)$}

La posición angular crece como rampa durante los pulsos de tensión, consecuencia directa del integrador puro ($\dot{\theta}_m = \omega_m$). Alcanza un máximo de $\approx 250$~rad en $t \approx 0.7$~s, que se mantiene hasta $t \approx 1.1$~s (durante $v_{qs}^* = 0$)  y retorna hacia cero durante el pulso negativo de $v_{qs}^*$. Ambos modelos presentan excelente concordancia. \textcolor{red}{La pequeña diferencia residual al final (offset constante una vez que $\theta_m$ regresa cerca de cero) se debe a la acumulación del error integrado entre las velocidades de ambos modelos, efecto inherente a la presencia del integrador puro en la planta.}

\begin{figure}[H]
    \centering
    \includegraphics[width=\textwidth]{imagenes/sim_theta_m.png}
    \caption{Posición angular $\theta_m(t)$ vs t: comparación modelo NL vs LTI aumentado ($i_{ds}^r(0) = 0$~A).}
    \label{fig:sim_theta_m}
\end{figure}

\textbf{Velocidad angular $\omega_m(t)$.}

La velocidad responde con dinámica dominada por la constante de tiempo mecánica $\tau_m = J_{eq}/b_{eq} \approx 0{,}70$~s. Durante el pulso positivo de $v_{qs}^*$ ($0{,}1$--$0{,}7$~s), $\omega_m$ crece hasta $\approx +400$~rad/s, modulada por las perturbaciones de torque de carga $T_{ld}$. Durante el pulso negativo ($1{,}1$--$1{,}7$~s), la respuesta es simétrica, alcanzando $\approx -400$~rad/s. La concordancia NL--LTI es excelente, confirmando la efectividad de la linealización por retroalimentación NL.

\begin{figure}[H]
    \centering
    \includegraphics[width=\textwidth]{imagenes/sim_omega_m.png}
    \caption{Velocidad angular $\omega_m(t)$ vs t: comparación modelo NL vs LTI aumentado ($i_{ds}^r(0) = 0$~A).}
    \label{fig:sim_omega_m}
\end{figure}

\textbf{Corriente $i_{qs}(t)$.}

La corriente $i_{qs}$ responde rápidamente a los escalones de $v_{qs}^*$ con constante de tiempo eléctrica $\tau_e = L_q/R_{s0} \approx 5{,}69$~ms, alcanzando su valor de régimen en $\approx 5\tau_e \approx 28{,}4$~ms. Se observan picos de $\approx \pm 10$~A. Ambos modelos se solapan casi perfectamente, confirmando que las leyes de control NL linealizan efectivamente el sistema. Dado que para $i_{ds}^r = 0$ el torque electromagnético se reduce a $T_{em} = k_t \cdot i_{qs}$ (con $k_t = \tfrac{3}{2} P_p \lambda_r$), la forma de onda de $T_{em}(t)$ es directamente proporcional a $i_{qs}(t)$.

\begin{figure}[H]
    \centering
    \includegraphics[width=\textwidth]{imagenes/sim_iqs.png}
    \caption{Corriente de eje $q$, $i_{qs}(t)$ vs t: comparación modelo NL vs LTI aumentado ($i_{ds}^r(0) = 0$~A).}
    \label{fig:sim_iqs}
\end{figure}

\textbf{Corrientes $i_{ds}(t)$ e $i_{0s}(t)$.}

Para la condición inicial $i_{ds}^r(0) = 0$~A, la corriente de eje $d$ permanece idénticamente nula en ambos modelos durante toda la simulación, ya que no existe excitación externa en el eje $d$ ($v_{ds}^* = 0$) y la ley de control mínima ($v_{ds} = -L_q i_{qs} \omega_r$) compensa exactamente el acoplamiento cruzado sin inyectar corriente en dicho eje. Análogamente, la corriente de secuencia cero $i_{0s}(t)$ permanece nula al no existir excitación homopolar ($v_{0s} = 0$) ni condición inicial no nula. El comportamiento de $i_{ds}^r(t)$ para condiciones iniciales $i_{ds}^r(0) \neq 0$ se analiza en detalle en el Sub-ítem~c.

\textbf{Temperatura del estator $T_s(t)$.}

La temperatura crece levemente respecto de $T_{amb} = 20\,^\circ$C debido a las pérdidas resistivas $P_{Cu} = \tfrac{3}{2}R_s(i_{qs}^2 + i_{ds}^{r\,2})$. El cambio es mínimo dado que la constante de tiempo térmica $\tau_{th} = R_{ts\text{-}amb} \cdot C_{ts} \approx 120$~s es muy superior a la duración de la simulación ($2{,}2$~s). Este resultado justifica la hipótesis del modelo LTI de utilizar $R_s = R_{sREF}$ constante: la variación de resistencia por efecto térmico es despreciable en esta escala temporal. Nótese que el modelo LTI aumentado (4 estados) \textbf{no incluye} la temperatura como variable de estado, ya que asume $R_s$ constante por diseño.

\begin{figure}[H]
    \centering
    \includegraphics[width=\textwidth]{imagenes/sim_Ts.png}
    \caption{Temperatura del estator $T_s(t)$ vs t: comparación modelo NL vs LTI aumentado ($i_{ds}^r(0) = 0$~A).}
    \label{fig:sim_Ts}
\end{figure}

\textbf{Tensión $v_{ds}(t)$ de la ley de control mínima.} 

La ley de control mínima impone internamente la tensión $v_{ds} = -L_q \, i_{qs} \, \omega_r$, proporcional al producto $i_{qs} \cdot \omega_m$. Esta tensión es no nula solo cuando ambas variables lo son simultáneamente, y su magnitud crece con el nivel de operación del sistema.

\begin{figure}[H]
    \centering
    \includegraphics[width=\textwidth]{imagenes/sim_v_ds.png}
    \caption{Tensión $v_{ds}$ vs t: modelo NL ($i_{ds}^r(0) = 0$~A).}
    \label{fig:sim_v_ds}
\end{figure}

\textbf{Corrientes en coordenadas $abc$.}

Mediante la transformada inversa de Park se obtienen las corrientes de fase reales que circularían por los bobinados del estator:
\begin{equation}
    \begin{bmatrix} i_a \\ i_b \\ i_c \end{bmatrix} = \mathbf{T}_{dq0}^{-1}(\theta_e) \begin{bmatrix} i_{ds}^r \\ i_{qs} \\ i_{0s} \end{bmatrix}, \quad \theta_e = P_p \cdot \theta_m
\end{equation}

\begin{figure}[H]
    \centering
    \includegraphics[width=\textwidth]{imagenes/sim_corrientes_abc.png}
    \caption{Corrientes de fase $i_a$, $i_b$, $i_c$ obtenidas por transformada inversa de Park (modelo NL, $i_{ds}^r(0) = 0$~A).}
    \label{fig:sim_corrientes_abc}
\end{figure}

Las formas de onda resultantes son sinusoidales moduladas en amplitud, con frecuencia eléctrica $\omega_e = P_p \cdot \omega_m$ proporcional a la velocidad del rotor. Cuando $\omega_m = 0$ (antes de $t = 0{,}1$~s y en los intervalos de velocidad nula), las corrientes $abc$ son constantes (frecuencia eléctrica nula). A medida que el motor acelera, la frecuencia de las oscilaciones crece proporcionalmente.

\textbf{Curva paramétrica: torque vs velocidad (cuadrantes de operación).}

La curva paramétrica $T_{em}(\omega_m)$ se construye a partir de los datos del script MATLAB, segmentando la trayectoria en 10 transitorios que conectan los puntos de (cuasi-)equilibrio $P_0$ a $P_{10}$.

\begin{figure}[H]
    \centering
    \includegraphics[width=\textwidth]{imagenes/sim_torque_vs_velocidad.png}
    \caption{Torque electromagnético vs velocidad angular: cuadrantes de operación del motor PMSM. Cada transitorio se distingue por color, con flechas indicando el sentido de recorrido y diamantes marcando los puntos de equilibrio $P_0$ a $P_{10}$.}
    \label{fig:sim_torque_vel}
\end{figure}

El sistema recorre los cuatro cuadrantes durante la secuencia completa de pulsos:
\begin{itemize}
    \item \textbf{Cuadrante I} ($+\omega_m, +T_{em}$): motorización en sentido positivo. El motor entrega torque en el mismo sentido que la velocidad.
    \item \textbf{Cuadrante II} ($-\omega_m, +T_{em}$): frenado regenerativo en sentido negativo.
    \item \textbf{Cuadrante III} ($-\omega_m, -T_{em}$): motorización en sentido negativo, análogo al Cuadrante~I para el pulso negativo de $v_{qs}^*$.
    \item \textbf{Cuadrante IV} ($+\omega_m, -T_{em}$): frenado regenerativo en sentido positivo. El torque frena la velocidad positiva residual.
\end{itemize}


\textbf{Ángulo de Carga $\delta(t)$}

El ángulo de carga $\delta(t)$ es el que representa al desfasaje electromagnético entre el ángulo eléctrico $\theta_{ev}$ del rotor respecto a la posicion angular del rotor en cordenadas eléctricas $\theta_r$. Donde se utiliza como referencia electrica, la posición del voltaje $V_{as}$, que es constante en estado estacionario.
\begin{equation}\label{eq:angulo_carga}
    \delta(t) \equiv \theta_{ev}(t)-\theta_r(t) = \int_0^t[\omega_e(\xi)-\omega_r(\xi)]d\xi +[\theta_e(0)-\theta_r(0)]
\end{equation}
Tal que:
\begin{itemize}
    \item $\omega_e$: velocidad angular de la referencia (tension fase \textbf{a})
    \item $\omega_r$: velocidad angular del rotor
\end{itemize}
Al mismo tiempo, se asumen posiciones angulares iniciales nulas para ambas variables.\\
\textbf{Cálculo de la coordenada eléctrica sincrónica $\theta_{ev}$:}
\begin{align*}
v_{as}(t) &= V_m \cos\big(\theta_{ev}(t)\big) \\
v_{bs}(t) &= V_m \cos\left(\theta_{ev}(t) - \frac{2\pi}{3}\right) \\
v_{cs}(t) &= V_m \cos\left(\theta_{ev}(t) + \frac{2\pi}{3}\right)
\end{align*}
Siendo $V_m$ el valor pico de la tensión senoidal aplicada.
Es necesario aplicar la Transformación de Clarke para proyectar las tensiones sobre un plano:
\begin{align*}
    V_{\alpha} = \frac{2}{3}\Big(v_{as} - \frac{1}{2}v_{bs}-\frac{1}{2}v_{cs}\Big) \\
    V_{\beta} = \frac{\sqrt{3}}{3}\big(v_{bs}-v_{cs}\big)
\end{align*}
Substituyendo obtenemos:
\begin{align}
    V_{\alpha} = V_m \cos{\theta_{ev}} \\
    V_{\beta} = V_m \sin{\theta_{ev}}
\end{align}
Por lo tanto para obtener el ángulo $\theta_{ev}$ basta con aplicar la funcion \textit{atan2} tal que:
\begin{equation}
    \theta_{ev} = \text{atan2}(V_{\beta},V_{\alpha})= \text{atan2}(V_m \sin{\theta_{ev}}, V_m \cos{\theta_{ev}})
\end{equation}

\begin{figure}[H]
    \centering
    \includegraphics[width=\textwidth]{imagenes/sim_angulo_carga.png}
    \caption{Ángulo de carga $\delta(t)$ vs t (modelo NL, $i_{ds}^r(0) = 0$~A).}
    \label{fig:sim_angulo_carga}
\end{figure}
% =====================================================================
\subsubsection{Métricas de transitorio (Sub-ítem b)}
\label{sec:sim_item_b}

Se determinan las métricas de los 10 transitorios para el caso nominal $i_{ds}^r(0) = 0$~A del modelo LTI aumentado. La Tabla~\ref{tab:metricas_trans} presenta, para cada evento de conmutación, el valor de $\omega_m$ al final del segmento temporal correspondiente, el valor de régimen teórico de $i_{qs}$, y las métricas de transitorio de $i_{qs}$ (tiempo de crecimiento $t_r$, tiempo de asentamiento $t_s$ y sobrepico $M_p$).

Dado que la dinámica mecánica es mucho más lenta que la eléctrica, $\omega_m$ no alcanza en la práctica su valor de régimen dentro de cada segmento de 200~ms ($\tau_m = J_{eq}/b_{eq} = 0{,}90$~s $\gg 0{,}2$~s), por lo que se reporta el valor instantáneo al final de cada intervalo. En contraste, $i_{qs}$ sigue una dinámica de primer orden con $\tau_e = L_q/R_{s0} = 5{,}69$~ms y alcanza la banda de $\pm 1\,\%$ en $t_s = \tau_e \ln(100) \approx 26{,}2$~ms, permitiendo medir métricas de transitorio estándar.

\begin{table}[H]
    \centering
    \caption{Métricas de transitorio para cada evento de conmutación (modelo LTI, $i_{ds}^r(0) = 0$~A).}
    \label{tab:metricas_trans}
    \renewcommand{\arraystretch}{1.15}
    \footnotesize
    \begin{tabular}{|c|l|r|r|r|r|r|}
        \hline
        \textbf{Trans.} & \textbf{Evento} & $\omega_m(t_f)$ & $i_{qs,ss}$ & $t_{r,i_{qs}}$ & $t_{s,i_{qs}}$ & $M_{p,i_{qs}}$ \\
         & & [rad/s] & [A] & [ms] & [ms] & [\%] \\
        \hline
        1 & $v_{qs}^* \to +19{,}6$~V   &  12220  & 19{,}21  & 12{,}5 & 26{,}2 & 0 \\
        2 & $T_{ld} \to +6{,}28$~Nm    &  21855  & 19{,}21  & --     & --     & -- \\
        3 & $T_{ld} \to -6{,}28$~Nm    &  30522  & 19{,}21  & --     & --     & -- \\
        4 & $v_{qs}^* \to 0$~V         &  25245  &  0{,}00  & 12{,}5 & 26{,}2 & 0 \\
        5 & $T_{ld} \to 0$~Nm          &  20222  &  0{,}00  & --     & --     & -- \\
        6 & $v_{qs}^* \to -19{,}6$~V   &   3979  & $-$19{,}21 & 12{,}5 & 26{,}2 & 0 \\
        7 & $T_{ld} \to +6{,}28$~Nm    &  $-$9828  & $-$19{,}21 & --     & --     & -- \\
        8 & $T_{ld} \to -6{,}28$~Nm    & $-$19939 & $-$19{,}21 & --     & --     & -- \\
        9 & $v_{qs}^* \to 0$~V         & $-$15818 &  0{,}00  & 12{,}5 & 26{,}2 & 0 \\
       10 & $T_{ld} \to 0$~Nm          & $-$11341 &  0{,}00  & --     & --     & -- \\
        \hline
    \end{tabular}

    \vspace{2mm}
    \scriptsize $t_f$: instante final de cada segmento. $i_{qs,ss}=v_{qs}^*/R_{s0}$: valor de régimen teórico del eje $q$. ``--'': métrica no aplicable porque el evento no modifica $v_{qs}^*$.
\end{table}

\textbf{Observaciones:}
\begin{itemize}
    \item La corriente $i_{qs}$ presenta respuesta de \textbf{primer orden pura} ($M_p = 0$\,\%) con $t_r = 12{,}5$~ms y $t_s = 26{,}2$~ms, consistentes con $\tau_e = L_q/R_{s0} = 5{,}69$~ms: en particular, $t_r \approx \tau_e \ln 9$ y $t_s(1\,\%) \approx \tau_e \ln 100$. Estos valores son \textbf{idénticos} para todos los escalones de $v_{qs}^*$ (transitorios 1, 4, 6, 9), confirmando la linealidad del sistema.

    \item Los eventos de $T_{ld}$ \textbf{no afectan} a $i_{qs}$ (métricas no aplicables), confirmando el desacoplamiento estructural: $T_{ld}$ entra por la ecuación mecánica y no tiene camino hacia la dinámica eléctrica del eje $q$.
\end{itemize}

\textbf{Influencia relativa de las acciones externas:}

\begin{itemize}
    \item \textbf{Tensión $v_{qs}^*(t)$}: actúa directamente sobre $i_{qs}$ con dinámica rápida ($\tau_e = 5{,}69$~ms). Es la acción \textbf{dominante}: genera $T_{em} = k_t \cdot i_{qs}$, que impulsa la dinámica mecánica a través del integrador.

    \item \textbf{Torque de carga $T_{ld}(t)$}: actúa directamente sobre $\omega_m$ a través de la ecuación mecánica ($\tau_m = J_{eq}/b_{eq} = 0{,}90$~s). Es una perturbación que modifica la trayectoria de velocidad sin alterar $i_{qs}$.

    \item La separación de escalas temporales ($\tau_m/\tau_e \approx 159$) explica por qué las corrientes responden casi instantáneamente a cambios de tensión, mientras que la velocidad presenta transitorios mucho más lentos.
\end{itemize}

% =====================================================================
\subsubsection{Efecto de $i_{ds}^r(0)$ sobre el sistema (Sub-ítem c)}
\label{sec:sim_item_c}

Se compara el comportamiento de $i_{ds}^r(t)$ para las tres condiciones iniciales en ambos modelos.

\begin{figure}[H]
    \centering
    \includegraphics[width=0.9\textwidth]{imagenes/sim_ids_NL_3casos.png}
    \caption{Comparación de $i_{ds}^r(t)$ para $i_{ds}^r(0) \in \{0,\; +0{,}5,\; -0{,}5\}$~A: modelo NL.}
    \label{fig:sim_comparacion_ids}
\end{figure}

\begin{figure}[H]
    \centering
    \includegraphics[width=0.9\textwidth]{imagenes/sim_ids_LTI_3casos.png}
    \caption{Comparación de $i_{ds}^r(t)$ para $i_{ds}^r(0) \in \{0,\; +0{,}5,\; -0{,}5\}$~A: modelo LTI.}
    \label{fig:sim_ids_zoom}
\end{figure}

\textbf{Análisis:}
\begin{itemize}
    \item En el \textbf{modelo LTI}, $i_{ds}^r(t)$ decae como exponencial pura:
    \begin{equation}
        i_{ds}^r(t) = i_{ds}^r(0) \cdot e^{-t/\tau_d}, \quad \tau_d = \frac{L_d}{R_{s0}} = \frac{6{,}6 \times 10^{-3}}{1{,}02} = 6{,}47 \text{ ms}
    \end{equation}
    completamente \textbf{desacoplado} del resto de los estados. Esto es consecuencia de la estructura diagonal por bloques de la matriz $A_{aug}$: el polo $p_d = -R_{s0}/L_d \approx -154{,}5$~rad/s es real, estable y no comparte autovector con ningún otro estado.

    \item En el \textbf{modelo NL}, $i_{ds}^r(t)$ también decae exponencialmente con la misma constante de tiempo. Aunque existe un acoplamiento residual NL del tipo $\omega_r \cdot L_q \cdot i_{qs}$, la ley de control mínima ($v_{ds} = -L_q i_{qs} \omega_r$) lo cancela exactamente, y la ley complementaria ($v_{qs} = v_{qs}^* + \omega_r L_d i_{ds}^r$) elimina el efecto cruzado inverso.

    \item En ambos modelos, el efecto de $i_{ds}^r(0) \neq 0$ se extingue en $\approx 5\tau_d \approx 32$~ms, sin impacto apreciable sobre la dinámica mecánica ($\omega_m$, $\theta_m$) ni sobre $i_{qs}$. Esto justifica la hipótesis adoptada en la sección de linealización: la condición $i_{ds}^r(0) = 0$ no es restrictiva en la práctica, ya que cualquier valor inicial decae rápidamente a cero.
\end{itemize}

% =====================================================================
\subsubsection{Field Forcing y Field Weakening (Sub-ítem d)}
\label{sec:sim_item_d}

Se agrega una consigna de tensión en eje $d$, sumada a la ley de control mínima:
\begin{equation}
    v_{ds}(t) = \underbrace{-L_q \, i_{qs} \, \omega_r}_{\text{ley mínima}} + v_{ds}^*(t), \quad
    v_{ds}^*(t) = \begin{cases}
    0 & t < 0{,}5 \text{ s} \\
    +1{,}9596 \text{ V} & t \geq 0{,}5 \text{ s} \quad \text{(field forcing)}\\
    -1{,}9596 \text{ V} & t \geq 0{,}5 \text{ s} \quad \text{(field weakening)}
    \end{cases}
\end{equation}

donde $V_{ds}^* = V_{qs,nom}/10 = 1{,}9596$~V. Se simula con $i_{ds}^r(0) = 0$~A para aislar el efecto de la tensión inyectada.

\begin{figure}[H]
    \centering
    \includegraphics[width=\textwidth]{imagenes/sim_ff_ids.png}
    \caption{Efecto del field forcing ($v_{ds}^* = +1{,}96$~V) y field weakening ($v_{ds}^* = -1{,}96$~V) sobre $i_{ds}^r(t)$, comparado con el caso base ($v_{ds}^* = 0$).}
    \label{fig:sim_field_forcing}
\end{figure}

\textbf{Análisis:}
\begin{itemize}
    \item En el \textbf{modelo LTI}, $v_{ds}^*$ afecta \textbf{únicamente} a $i_{ds}^r$, que crece exponencialmente hasta su valor estacionario $i_{ds,ss}^r = v_{ds}^*/R_{s0} = \pm 1{,}92$~A. Los demás estados ($\theta_m$, $\omega_m$, $i_{qs}$) permanecen \textbf{inalterados} por el desacoplamiento estructural de la matriz $A_{aug}$.

    \item En el \textbf{modelo NL}, $i_{ds}^r \neq 0$ introduce acoplamientos cruzados que el LTI no captura:
    \begin{itemize}
        \item \textit{Torque de reluctancia}: $\Delta T_m = \tfrac{3}{2}P_p(L_d - L_q) i_{ds}^r i_{qs}$, que modifica el torque electromagnético neto.
        \item \textit{FEM de acoplamiento}: los términos $\omega_r L_d i_{ds}^r$ en eje $q$ y $-\omega_r L_q i_{qs}$ en eje $d$ generan tensiones inducidas cruzadas.
    \end{itemize}

    \item Para \textbf{field forcing} ($v_{ds}^* > 0$): $i_{ds}^r > 0$ genera un torque de reluctancia positivo (dado que $L_d > L_q$) que incrementa el torque total, produciendo sobrepicos mayores en $\omega_m$ y un régimen estacionario levemente diferente.

    \item Para \textbf{field weakening} ($v_{ds}^* < 0$): el efecto es simétrico opuesto. La corriente $i_{ds}^r < 0$ debilita el campo magnético efectivo, reduciendo el torque disponible.

    \item Ambos efectos son pequeños para la amplitud utilizada ($V_{qs,nom}/10$), pero la diferencia entre modelos demuestra las \textbf{limitaciones del LTI}: al operar con $i_{ds}^r$ alejado de cero, los términos no lineales omitidos por la linealización se vuelven relevantes.
\end{itemize}

%%%SECCION 516


\section{Diseño, Análisis y Simulación de Controlador de Movimiento en Cascada con Modulador de Torque}
\subsection{Diseño del Modulador de Torque}

En esta sección se desarrolla el diseño e implementación del Modulador de Torque Equivalente, el cual operará como el controlador interno vectorial de corriente y torque del sistema. El objetivo principal de esta etapa es recibir una consigna de torque de referencia $(T_m^*)$ y generar las tensiones de comando necesarias en coordenadas rotóricas qd0 del estator, logrando un lazo de control rápido y libre de acoplamientos indeseados. Para lograrlo, se implementa una estrategia de desacoplamiento completo que compensa activamente todas las retroalimentaciones físicas naturales de la planta, cancelando los acoplamientos magnéticos cruzados, las fuerzas contraelectromotrices y las caídas resistivas variables con la temperatura $(R_s(T_s^\circ))$
\subsection{Desacoplamiento de retroalimentaciones naturales}
Recordando las ecuaciones \eqref{eq:vqs_rotor}, \eqref{eq:vds_rotor} y \eqref{eq:v0s_rotor}
    \begin{equation*}
        v^r_{qs}=R_s(T_s^\circ(t))i^r_{qs}(t) + L_q\frac{d i^r_{qs}(t)}{dt}+[\lambda'^r_m + L_d i^r_{ds}(t)]\omega_r(t)
    \end{equation*}
    \begin{equation*}
        v_{ds}^r=R_s(T_s^\circ(t))i^r_{ds}(t)+L_d\frac{d i^r_{ds}(t)}{dt} - L_qi^r_{qs}(t)\omega_r(t)
    \end{equation*}
    \begin{equation*}
        v^r_{0s}=R_s(T_s^\circ(t))i^r_{0s}(t)+L_{ls}\frac{d i^r_{0s}(t)}{dt}
    \end{equation*}
Podemos notar las retroalimentaciones naturales, vease la dependencia de la Resistencia variable con la temperatura $T_s^\circ(t)$ y el acoplamiento entre las corrientes $i_{qs}$ y $i_{ds}$ y demás causales al mismo tiempo de no linealidades ya estudiadas anteriormente en el presente trabajo.

Por lo que la propuesta principal es cancelar estos términos que poseen retroalimentaciones físicas naturales y obtener ecuaciones que relacionen a la dinámica de las corrientes directamente con el voltaje, para (como a posteriori se verá) implementar el modulador de torque de forma sencilla.

Lo que nos lleva a proponer las siguientes ecuaciones, que reemplazandolas luego en \eqref{eq:vqs_rotor}, \eqref{eq:vds_rotor} y \eqref{eq:v0s_rotor} se obtendrán ecuaciones diferenciales en una forma sumamente útil.
\begin{equation}
v_{qs}^{r}(t) = v_{qs}^{r*}(t) 
+ R_s\!\big(T_s^{\circ}(t)\big)\, i_{qs}^{r}(t) 
+ \left[\lambda'_m + L_d\, i_{ds}^{r*}(t)\right] P_p\, \omega_m(t)
\label{eq:vqs_r}
\end{equation}

\begin{equation}
v_{ds}^{r}(t) = v_{ds}^{r*}(t) 
+ R_s\!\big(T_s^{\circ}(t)\big)\, i_{ds}^{r}(t) 
- L_q\, i_{qs}^{r*}(t)\, P_p\, \omega_m(t)
\label{eq:vds_r}
\end{equation}

\begin{equation}
v_{0s}(t) = v_{0s}^{*}(t) 
+ R_s\!\big(T_s^{\circ}(t)\big)\, i_{0s}(t)
\label{eq:v0s}
\end{equation}
Obteniendo si reemplazamos:
\begin{equation}
    v_{qs}^{r*}(t) = L_q \frac{di_{qs}^r(t)}{dt}
    \label{eq:vqs_ref}
\end{equation}
\begin{equation}
    v_{ds}^{r*}(t) = L_d \frac{di_{ds}^r(t)}{dt}
    \label{eq:vds_ref}
\end{equation}
\begin{equation}
    v_{0s}^{r*}(t) = L_{ls} \frac{di_{0s}^r(t)}{dt}
    \label{eq:v0s_ref}
\end{equation}
\begin{figure}[H]
    \centering
    \includegraphics[width=0.8\textwidth]{imagenes/desacop_realimentacion_fisicas_naturales.png}
\end{figure}
\begin{figure}[H]
    \centering
    \includegraphics[width=0.4\textwidth]{imagenes/desacop_realimentacion_fisicas_naturales_bloque.png}
\end{figure}
\subsubsection{Comparación: desacoplamiento de retroalimentaciones físicas naturales vs. linealización por retroalimentación NL}

Al contrastar la estrategia de linealización mediante el Modulador de Torque (que compensa las retroalimentaciones físicas naturales del sistema) con la linealización por retroalimentación NL mínima desarrollada en la etapa inicial a lazo abierto, aparecen diferencias relevantes en alcance matemático, flexibilidad operativa y manejo de la dinámica interna:

\begin{itemize}
    \item \textbf{Generalización del espacio de operación (liberación de $i_{ds}^r$):} 
    La estrategia de linealización previa imponía de forma estricta la condición $i_{ds}^r(t)=0$ para lograr el desacoplamiento entre el canal de flujo magnético y el de torque. En cambio, el desacoplamiento completo implementado en el Modulador de Torque generaliza el modelo y permite que la corriente del eje directo tome valores distintos de cero. Esto es clave porque habilita, de manera física, operaciones de debilitamiento o reforzamiento de campo según los requerimientos del accionamiento.
    
    \item \textbf{Grado de cancelación de las dinámicas naturales:} 
    La linealización por retroalimentación NL original se limitaba a cancelar los acoplamientos cruzados inductivos entre los ejes $q$ y $d$. El nuevo enfoque de desacoplamiento es más exhaustivo: inyecta tensiones de referencia calculadas para cancelar todas las retroalimentaciones naturales de la planta. Esto incluye no solo los acoplamientos cruzados, sino también la fuerza contraelectromotriz de los imanes permanentes y las caídas de tensión resistivas de los devanados.
    
    \item \textbf{Control explícito de la dinámica de corriente (asignación de polos):} 
    En la linealización inicial, la cancelación se realizaba sin ubicar explícitamente los polos del lazo interno de corriente, lo que limitaba el ajuste de la rapidez de respuesta. Con el desacoplamiento completo actual, las corrientes del estator pasan a responder de manera casi independiente, comportándose aproximadamente como integradores puros ante las entradas de tensión. Esto permite introducir un Modulador de Corriente estrictamente proporcional que ubica los polos del lazo interno en posiciones específicas (por ejemplo, $p_i=-5000$ rad/s), otorgando mayor control sobre la velocidad de respuesta en transitorios.
\end{itemize}


\subsubsection{Efecto de estimar la dependencia $R_s(T_s(t))$ vs. utilizar un valor nominal constante}

Al diseñar el Modulador de Torque Equivalente encargado del desacoplamiento, surge la disyuntiva de cómo tratar la resistencia del estator. En la planta física, la resistencia $R_s$ depende de forma inevitable de la temperatura $T_s(t)$ debido al calentamiento progresivo por efecto Joule. El tratamiento de este parámetro define la robustez del controlador:

\begin{itemize}
    \item \textbf{Estimación en línea (resistencia variable):} 
    Si el controlador interno considera esta dependencia térmica y actualiza en tiempo real el valor de $R_s(T_s(t))$ a partir de las mediciones del sensor de temperatura, el cálculo de las tensiones de desacoplamiento resulta matemáticamente más preciso. Esto deriva en un comportamiento más robusto frente a cambios térmicos, al reducir perturbaciones cruzadas y errores de seguimiento durante transitorios y en régimen permanente.
    
    \item \textbf{Uso de valor nominal (resistencia constante):} 
    Si el controlador asume un valor nominal constante (típicamente el medido a temperatura ambiente), la compensación interna pierde exactitud a medida que la temperatura del devanado aumenta y se aleja del punto de referencia. En consecuencia, la tensión calculada para compensar la caída resistiva resulta insuficiente frente a la caída de tensión real en una máquina caliente.
\end{itemize}
%insertar aqui graficas
\begin{figure}[H]
    \centering
    \includegraphics[width=0.6\textwidth]{imagenes/comp_torque_rs.png}
    \caption{Evolución temporal del par neto ($T_{\text{neto}}$) en [N.m]. Se contrasta el comportamiento bajo condiciones de resistencia estatórica constante (RS Cte., línea amarilla) y resistencia estatórica variable (RS Variable, línea azul).}
    \label{fig:comparacion_torque_neto_RS}
\end{figure}

\begin{figure}[H]
    \centering
    \includegraphics[width=0.6\textwidth]{imagenes/comp_velocidad_rs.png}
    \caption{Comportamiento de la velocidad angular mecánica ($\omega_m(t)$) en [rad/s] en función del tiempo. Se comparan las respuestas del sistema operando con resistencia estatórica constante (RS Cte., línea azul) y con resistencia estatórica variable (RS Variable, línea roja).}
    \label{fig:comparacion_velocidad_angular_RS}
\end{figure}
El impacto de no compensar la variación térmica se manifiesta en una degradación del desempeño global del sistema. Esta inexactitud en la cancelación resistiva produce un error en estado estacionario en las corrientes internas y, consecuentemente, en el torque electromagnético generado.

Como puede verse en la Fig. \ref{fig:comparacion_torque_neto_RS} y \ref{fig:comparacion_velocidad_angular_RS} este déficit sostenido de torque aplicado tiene una repercusión directa y acumulativa sobre el comportamiento mecánico. Al entregar un par ligeramente inferior al demandado por la consigna, se genera una divergencia progresiva en las curvas de velocidad angular. La trayectoria de velocidad del sistema que opera con un controlador de resistencia constante se desvía de forma creciente respecto de la trayectoria ideal, lo que evidencia la necesidad de un modelado térmico preciso para garantizar el desempeño exigido en regímenes de operación prolongados.
\subsection{Diseño de los Lazos de Control de Corrientes $\bm{i}_{qd0s}^r$}
Dado que nuestra planta es de tipo 1 es posible aplicar una simple Ley de Control proporcional para lograr un error de estado estacionario nulo, por lo que podemos plantear:
\begin{equation}
    \bm{v}_{qd0s}^{r*}(t) = \bm{K_p} \big(\bm{i}_{qd0s}^{r*}(t)-\bm{i}_{qd0s}^{r}(t)\big)
\end{equation}
En forma escalar:
\begin{equation*}
\begin{aligned}
v_{qs}^{r*}(t) &= K_{pq} \left(i_{qs}^{r*}(t) - i_{qs}^{r}(t)\right) \\
v_{ds}^{r*}(t) &= K_{pd} \left(i_{ds}^{r*}(t) - i_{ds}^{r}(t)\right) \\
v_{0s}^{r*}(t) &= K_{p0} \left(i_{0s}^{r*}(t) - i_{0s}^{r}(t)\right)
\end{aligned}
\end{equation*}
Lo que nos lleva a determinar la siguiente funcion de transferencia, reemplazando las ecuaciones \eqref{eq:vqs_ref}, \eqref{eq:vds_ref}, \eqref{eq:v0s_ref} en la Ley de Control:
\begin{gather*}
    L_qsI_{qs}^r(s) = K_{pq}\big(I_{qs}^{r*}(s)-I_{qs}^{r}(s)\big) \\ 
    \Big(L_qs+K_{pq}\Big)I_{qs}^r(s)=K_{pq}I_{qs}^{r*}(s) \\
    \frac{I_{qs}^r(s)}{I_{qs}^{r*}(s)} = \frac{1}{1+\frac{L_q}{K_{pq}}s}
\end{gather*}
Idem para los demas ejes.
Queremos los polos en $p_{1,2,3} = -5000 \ [\frac{rad}{s}]$, por lo tanto las ganancias $K_{p}^{qd0}$ quedarían: \[K_{p}^{qd0}=L_{q,d,ls}\cdot 5000 \ \Omega\]
\begin{figure}[H]
    \centering
    \includegraphics[width=0.6\textwidth]{imagenes/mod_corriente.png}
\end{figure}
\begin{figure}[H]
    \centering
    \includegraphics[width=0.2\textwidth]{imagenes/mod_corriente_subsystem.png}
\end{figure}
\subsection{Incorporación de Consigna de Torque: Modulador de Torque}
Para el sistema electromagnético la unica forma de variar el torque es a traves de la corriente, por lo tanto para incorporar esta nueva variable manipulada, es necesario diseñar un simple Modulador de Torque, que traducirá nuestra consigna a una consigna de corriente $i_{qs}^{r*}$. Se propone entonces, aplicar una consigna de \textit{Torque Mecánico}, para que, al momento de aplicar la consigna, no sea necesario tener en cuenta la fricción y la gravedad, por lo que implementamos un \textit{precomputed torque}.
\begin{equation}
    T_m^{\ast}(t) = T_m^{\ast\prime}(t) + b_{eq}\omega_m(t) + \frac{g\cdot k_l}{r}\sin{\Bigg( \frac{\theta_m(t)}{r}\Bigg)}
\end{equation}
Como se dijo, la idea principal es obtener una consigna de corriente, donde, aplicando Control Vectorial, se conseguirá una consigna de corriente en cuadratura $i_{qs}^{r*}$. Si despejamos de la ecuación \eqref{eq:torque_electromagnetico}.
    \begin{equation*}
            T_m^*(t) = \frac{3}{2} P_p \cdot \lambda'^r_m \cdot i^r_{qs}(t) + \frac{3}{2} P_p (L_d-L_q)\cdot i^r_{ds}(t)\cdot i^{r*}_{qs}(t)
            \label{eq:torque_electromagnetico}
    \end{equation*}
\begin{gather*}
i_{qs}^{r*}(t)=
\frac{T_m^*(t)}
{\frac{3}{2}P_p\left[\lambda_m+(L_d-L_q)i_{ds}^{r}(t)\right]} \\
i_{qs}^{r*}(t)= \frac{T_m^{\ast\prime}(t) + b_{eq}\omega_m(t) + \frac{g\cdot k_l}{r}\sin{\Big( \frac{\theta_m(t)}{r}\Big)}}{\frac{3}{2}P_p\left[\lambda_m+(L_d-L_q)i_{ds}^{r}(t)\right]}
\end{gather*}
Puede analizarse en la ecuación anterior, que la variable $i_{ds}$ quedaría como un parámetro más, lo que habilita a la posiblidad de 2 casos ya vistos anteriormente en el presente trabajo.
\begin{itemize}
    \item $i_{ds} = 0$: Si se aplica esta consigna entonces el torque depende unicamente de la corriente en eje q y del flujo de los imanes permanentes, siendo suficiente para la mayoria de casos
    \item $i_{ds} \neq 0$: Ya sea por reforzamiento o debilitamiento de campo, $i_{ds}$ haria de una consigna más, para situaciones en las que el rango de voltaje del inversor no sea suficientemente amplio.
\end{itemize}
\begin{figure}[H]
    \centering
    \includegraphics[width=0.75\linewidth]{imagenes/mod_torque_diagram.png}
    \caption{Diagrama de bloques del modulador de torque implementado}
\end{figure}
% \begin{gather*}
%     J_{eq}\ddot\theta_m(t) = T_m^*(t) - \frac{T_l(t)}{r} \\
%     J_{eq}s^2\Theta_m(s)=T_m^*(s) - \frac{T_l(s)}{r} \\
%     J_{eq}s^2\Theta_m(s) =e_\theta(s)(b_as + K_{sa}+K_{sia}s^{-1})- \frac{T_l(s)}{r}\\
%     J_{eq}s^3\Theta_m(s) =e_\theta(s)(b_as^2 + K_{sa}s+K_{sia})- \frac{T_l(s)s}{r}\\
%     J_{eq}s^3\Theta_m(s) = [\Theta_m^{\ast}(s)-\Theta_m(s)](b_as^2 + K_{sa}s+K_{sia})- \frac{T_l(s)s}{r}\\
%     \Theta_m(s)\Big(J_{eq}s^3+b_as^2 + K_{sa}s+K_{sia}\Big)=\Theta_m^{\ast}(s)\Big(b_as^2 + K_{sa}s+K_{sia}\Big)- \frac{T_l(s)s}{r} \\
%     \Theta_m(s)=\frac{\Big(b_as^2 + K_{sa}s+K_{sia}\Big)}{J_{eq}s^3+b_as^2 + K_{sa}s+K_{sia}}\Theta_m^{\ast}(s)+\frac{-s/r}{J_{eq}s^3+b_as^2 + K_{sa}s+K_{sia}}T_l(s)
% \end{gather*}
% Lo que nos evidencia ambas funciones de transferencia, por lo tanto, nos podemos quedar con la de la entrada:
% \begin{equation}
%     G_{\theta_{m}^\ast} = \frac{\Theta_m(s)}{\Theta_m^\ast(s)} = \frac{\Big(b_as^2 + K_{sa}s+K_{sia}\Big)}{J_{eq}s^3+b_as^2 + K_{sa}s+K_{sia}}
% \end{equation}
% Si tenemos en cuenta los requerimientos de la consigna y expresamos el polinomio ideal basándonos en $\zeta$ y $\omega_n$:
% \begin{equation}
%     p_{ideal}(s)= (s+\omega_n)(s^2+2\zeta\omega_ns+\omega_n)=s^3+(2\zeta+1)\omega_ns^2+(2\zeta+1)\omega_n^2s+\omega_n^3
% \end{equation}
% Por lo que comparándolo con el polinomio mónico de $G_{\theta_{m}^\ast}$, lo que buscamos es: 
% \begin{gather*}
%     s^3+b_a/J_{eq}s^2+K_{sa}/J_{eq}s+K_{sia}/J_{eq}= s^3+(2\zeta+1)\omega_ns^2+(2\zeta+1)\omega_n^2s+\omega_n^3 \\
%     b_a = J_{eq}(2\zeta + 1)\omega_n \\
%     K_{sa}=J_{eq}(2\zeta+1)\omega_n^2 \\ 
%     K_{sia} = J_{eq}\omega_n^3
% \end{gather*}
\subsection{Controlador externo de posición y velocidad}

Dado que el objetivo del presente trabajo es controlar la posición y la velocidad del eje, se propone la implementación de un controlador PID mediante el método de sintonía serie. El objetivo es que dicho controlador manipule directamente el torque, aplicando consignas sobre el \textit{Modulador de Torque}. Para ello, se propone la siguiente ley de control:
\begin{equation}
    T_m^{*\prime}(s)= b_a\,e_\omega(s) + K_{sa}\,e_{\theta_m}(s)+K_{sia}\,\frac{e_{\theta_m}(s)}{s}
\end{equation}

En la consigna se especifica que el controlador debe presentar un amortiguamiento $\zeta = 0.75$ y una frecuencia natural $\omega_n = 800\ \frac{rad}{s}$. Para obtener tales características, puede emplearse el método de separación de anchos de banda, también denominado ``sintonía serie''.

Este método plantea una anidación entre los lazos de control. En este caso, se definen tres frecuencias significativas para el controlador:
\begin{equation}\label{eq:frecuencias_pid}
\omega_{vel}:= \frac{b_a}{J_{eq}}, \qquad
\omega_{pos}:= \frac{K_{sa}}{b_a}, \qquad
\omega_{int}:= \frac{K_{sia}}{K_{sa}}
\end{equation}

En este esquema, el parámetro $n$ indica cuán rápidos son los lazos internos en comparación con el lazo de posición. Así, el lazo de velocidad debe ser $n$ veces más rápido que el lazo de posición, mientras que el lazo integral debe ser $n$ veces más lento que este último. Por lo tanto, se adopta:
\begin{gather}
    \omega_{vel} := n\,\omega_{pos} \\
    \omega_{int} := \frac{1}{n}\,\omega_{pos}
\end{gather}

Dado que el desempeño especificado corresponde al lazo de posición, se toma:
\begin{equation}
    \omega_{pos} = \omega_n = 800\ \frac{rad}{s}
\end{equation}

Reemplazando estas relaciones en \eqref{eq:frecuencias_pid}, se obtienen las ganancias del controlador en función de $n$:
\begin{gather}
    b_a = J_{eq}\,\omega_{vel} = J_{eq}\,n\,\omega_{pos}
    \label{eq:ba_n} \\
    K_{sa} = b_a\,\omega_{pos} = J_{eq}\,n\,\omega_{pos}^2
    \label{eq:ksa_n} \\
    K_{sia} = K_{sa}\,\omega_{int} = J_{eq}\,\omega_{pos}^3
    \label{eq:ksia_n}
\end{gather}

Puede observarse que $b_a$ y $K_{sa}$ dependen del parámetro $n$, mientras que $K_{sia}$ queda determinado únicamente por $J_{eq}$ y por la frecuencia impuesta al lazo de posición.

Para justificar el valor de $n$, se considera el polinomio característico del controlador:
\begin{equation}
    p(s)=J_{eq}s^3+b_as^2+K_{sa}s+K_{sia}
\end{equation}

Reemplazando \eqref{eq:ba_n}, \eqref{eq:ksa_n} y \eqref{eq:ksia_n}, resulta:
\begin{equation}
    p(s)=J_{eq}\left(s^3+n\omega_{pos}s^2+n\omega_{pos}^2s+\omega_{pos}^3\right)
\end{equation}

Este polinomio puede factorizarse como:
\begin{equation}
    p(s)=J_{eq}(s+\omega_{pos})\left(s^2+(n-1)\omega_{pos}s+\omega_{pos}^2\right)
\end{equation}

Comparando el factor cuadrático con la forma estándar de un sistema de segundo orden:
\begin{equation}
    s^2+2\zeta\omega_ns+\omega_n^2
\end{equation}

y considerando que $\omega_n=\omega_{pos}$, se obtiene:
\begin{equation}
    (n-1)\omega_{pos}=2\zeta\omega_{pos}
\end{equation}

de donde resulta:
\begin{equation}
    n=2\zeta+1
\end{equation}

Finalmente, reemplazando el amortiguamiento especificado:
\begin{equation}
    n=2\cdot0.75+1=2.5
\end{equation}

Por lo tanto, para el caso nominal, las ganancias del controlador resultan:
\begin{equation*}
\begin{aligned}
    b_a &= 0.0396\ \frac{N\cdot m}{rad/s} \\
    K_{sa} &= 31.656\ \frac{N\cdot m}{rad} \\
    K_{sia} &= 10129.778\ \frac{N\cdot m}{rad\cdot s}
\end{aligned}
\end{equation*}



\subsection{Incorporación y diseño de Observador de Estado de Orden reducido para $\theta_m(t)$ y $\omega_n(t)$}
Se pretende estimar la velocidad del motor a partir del sensado de la posición, suponiendo por lo tanto, que no se dispone de un tacómetro. Esta estimación es posible debido al análisis que se realizó anteriormente (\textit{Controlabilidad y Observabilidad}) donde también se demuestra que no es observable la posición a partir de la velocidad.
\subsubsection{Modelo del subsistema mecánico}
Para implementar un observador de orden reducido no es necesario plantear todas las ecuaciones de estado, en este caso, el único subsistema involucrado en la estimación es el mecánico, que, si consideramos que ya se desacopló el término $-b_{eq}\omega_n(t)$ nos queda:
\begin{gather}
\begin{aligned}
    \dot\theta_m(t) &= \omega_n(t) \\
    \dot \omega_m(t) &= \frac{1}{J_{eq}}\Big(T'_m(t)-\frac{T_l(t)}{r}\Big)
\end{aligned}
\end{gather}
Que en forma matricial:
\begin{gather*}
    \bm{\dot x}
=
\begin{bmatrix}
    0 & 1 \\ 0 & 0
\end{bmatrix}
\bm{x}(t)
+
\begin{bmatrix}
    0 \\ \frac{1}{J_{eq}} 
\end{bmatrix}
    u(t)
\\
y(t) =
\begin{bmatrix}
    1 & 0
\end{bmatrix}
\bm{x}(t)
\end{gather*}
Donde:
\begin{equation*}
\bm{\dot x}(t)=
\begin{bmatrix}
    \dot\theta_m(t) \\\dot\omega_m(t)
\end{bmatrix}    ;\ \ 
\bm{x}(t)=
\begin{bmatrix}
    \theta_m(t) \\\omega_m(t)
\end{bmatrix}; \ \
u(t) = T_m'(t); \ \
\bm{x}(0) = \begin{bmatrix}
    \theta_m(0) \\\omega_m(0)
\end{bmatrix};\ \
y(t) = \theta_m(t);\ \  
A=\begin{bmatrix}
    0 & 1\\ 0 &0
\end{bmatrix}; \ \
C=\begin{bmatrix}
    1 &0
\end{bmatrix}
\end{equation*}
    
\subsubsection{Diseño del observador}
    El observador de estados reducido se basará en el observador de Luenberger, cuya dinámica sigue una Ley de Control proporcional al error entre la salida medida y la estimada.
    \begin{gather*}
    \begin{aligned}
    \dot{\bm{\bar x}}&=
    \mathrm{\bm{A}}
    \bm{\bar x} + 
    \begin{bmatrix}
    0 \\ \frac{1}{J_{eq}} 
    \end{bmatrix} u(t)+
    \mathrm{\bm{L}}\big(y -
    \mathrm{\bm{C}}\bm{\bar x}\big)\\
        \dot{\bm{\bar x}}&=
    \big(\mathrm{\bm{A}-\mathrm{\bm{L}}\mathrm{\bm{C}}\big)}\bm{\bar x} + 
    \begin{bmatrix}
    0 \\ \frac{1}{J_{eq}} 
    \end{bmatrix} u(t)+\mathrm{\bm{L}}\cdot y
    \end{aligned}
    \end{gather*}
    Definiendo a $\mathrm{\bm{L}}$ como la ganancia proporcional:
    \begin{equation*}
        \mathrm{\bm{L}}=
        \begin{bmatrix}
            K_{e\theta_m}\\
            K_{e\omega_m}
        \end{bmatrix}
    \end{equation*}
    Lo que resulta en el siguiente polinomio característico:
    \begin{gather*}
    \begin{aligned}
        \det \Big(\mathrm{s\bm{I-\big(A}}-\mathrm{\bm{LC}\big)}\Big)&=0,\\
        \begin{bmatrix}
            s & 0 \\ 
            0 & s
        \end{bmatrix}-
        \begin{bmatrix}
            0 & 1 \\
            0 & 0
        \end{bmatrix}+
        \begin{bmatrix}
            K_{e\theta_m} & 0 \\
            K_{e\omega_m} & 0
        \end{bmatrix} &=0,
        \\
        s^2 + s K_{e_{\theta_m}}+K_{e\omega_m}&=0
    \end{aligned}
    \end{gather*}
    Se propone en la consigna, sintonizar estas ganancias de tal manera de obtener ambos polos en $p_{1,2}=-3200 \frac{rad}{s}$, y siendo el polinomio deseado:
    \begin{gather*}
        p_{des} = (s-p_{1,2})^2=s^2-2p_{1,2}s+p_{1,2}^2
    \end{gather*}
    Comparando con el pol. característico:
    \begin{equation}
        s^2 + s K_{e_{\theta_m}}+K_{e\omega_m}=s^2-2p_{1,2}s+p_{1,2}^2
    \end{equation}
    Resultando entonces en las siguientes ganancias:
    \begin{gather}
        K_{e\theta_m} = 6400 \ \frac{rad}{s}\\K_{e\omega_m}=3200^2 \ \frac{rad}{s}
    \end{gather}
    \begin{figure}[H]
        \centering
        \includegraphics[width=0.8\linewidth]{imagenes/observador_diagram.png}
        \caption{Diagrama de bloques del observador de estados reducido}
    \end{figure}
    \subsection{Simulación en tiempo continuo con modelo completo NL}
    \subsubsection{Consigna trapezoidal de posición}
    \begin{figure}[H]
        \centering
        \includegraphics[width=1\linewidth]{imagenes/pos-velocidad-simulacion_cont.png}
        \caption{Posición y velocidad de la articulación en función del tiempo siguiendo la consigna trapezoidal propuesta}
    \end{figure}
    \begin{figure}[H]
        \centering
        \includegraphics[width=1\linewidth]{imagenes/overshoot_transitorio_trap_pos.png}
        \caption{Overshoot del tansitorio rampa-meseta de la posición de la articulación}
    \end{figure}
    \begin{figure}[H]
        \centering
        \includegraphics[width=1\linewidth]{imagenes/torque_pos_relacion.png}
        \caption{Torque neto aplicado por la articulación para seguir la consigna de posición.}
    \end{figure}
Podemos concluir a traves de las distintas gráficas que el sistema sigue la consigna de forma esperada. Aun así, todavia quedan problemas que solucionar, como por ejemplo, los picos de corriente que se producen en los transitorios.
\begin{figure}[H]
    \centering
    \includegraphics[width=1\linewidth]{imagenes/picos_corriente_trapz_pos.png}
    \caption{Picos de corriente rms en transitorios}
\end{figure}
\subsubsection{Rechazo a perturbaciones $T_{ld}$ en escalón}
Se eligió una combinación de escalones para hacer mas evidente el efecto del torque de perturbación y la \textit{performance} del controlador al seguir la consgina de posición.
\[
T_{ld}(t)=
\begin{cases}
0 \ \ \text{[N.m]}, & 0 \le t < 1,\\[4pt]
5 \ \ \text{[N.m]}, & 1 \le t < 2,\\[4pt]
2.5 \ \ \text{[N.m]}, & 2 \le t < 3,\\[4pt]
-2.5 \ \ \text{[N.m]}, & 3 \le t < 5,\\[4pt]
-5 \ \ \text{[N.m]}, & 5 \le t < 6,\\[4pt]
2.5 \ \ \text{[N.m]}, & 6 \le t \le 15.
\end{cases}
\]

\begin{figure}[H]
    \centering
    \includegraphics[width=0.75\linewidth]{imagenes/plot_perturb_error.png}
    \caption{Influencia del torque de perturbación en el seguimiento de la consigna de posición y el error creciente}
\end{figure}
\begin{figure}[H]
    \centering
    \includegraphics[width=0.75\linewidth]{imagenes/perturb_y_posicion_plot.png}
    \caption{{Perturbación, consigna y posición medida, puede verse claramente el error que aumenta con el tiempo}}
\end{figure}

\section{Verificación de desempeño y especificaciones de operación}
La primer variable que puede apreciarse que supera los límites de operación es la corriente, para limitar este estado del sistema lo primero que se implementará es una saturación en el voltaje que se requiere del inversor, para que este no supere la especificación debe imponerse la siguiente condición:
\begin{equation*}
\begin{aligned}
    |v_{as}(t)|,|v_{bs}(t)|,|v_{cs}(t)|\le \frac{\sqrt{2}\cdot V_{sl,max}}{\sqrt{3}}\\
    |v_{as}(t)|,|v_{bs}(t)|,|v_{cs}(t)|\le 39.19 \ V
    \end{aligned}
\end{equation*}
\begin{figure}[H]
    \centering
    \includegraphics[width=1\linewidth]{imagenes/corrientes_limitando_voltaje.png}
    \caption{Comparación de la corriente rms del estator limitando el voltaje en la Run 20 y sin limitar el voltaje en la Run 17.}
    \label{fig:corrientes_voltaje}
\end{figure}
\begin{figure}[H]
    \centering
    \includegraphics[width=1\linewidth]{imagenes/pos_voltaje_saturado.png}
\end{figure}
\begin{figure}[H]
    \centering
    \includegraphics[width=1\linewidth]{imagenes/pos_voltaje_nosaturado.png}
\end{figure}
Como puede verse en la \ref{fig:corrientes_voltaje} aunque se reduce bastante el tamaño del pico de corriente y además no cambia el error de posición entre la simulación con voltaje saturado, aun así no es suficiente para cumplir con las especificaciones de corriente.
Posteriormente se propone saturar la consigna de corriente en el modulador de torque, para que al menos no se imponga una corriente que desde un principio supere la corriente pico del motor.
\begin{figure}[H]
    \centering
    \includegraphics[width=1\linewidth]{imagenes/corrientes_comparacion.png}
    \caption{Puede verse que aun con una saturación en la consigna de corriente, le es imposible al accionamiento cumplir con la consigna sin superar sus especificaciones}
\end{figure}
Finalmente lo que nos lleva a atacar la forma de la consigna, si analizamos las gráficas previas en la simulación, puede verse como se le pide por momentos aceleraciones infinitas al accionamiento (saltos discontinuos en la velocidad). 

Inicialmente se propuso implementar perfil trapezoidal de la velocidad, obteniendo resultados sustancialmente superiores en términos de cumplimiento de especificaciones.
\begin{figure}[H]
    \centering
    \includegraphics[width=0.75\linewidth]{imagenes/perfil_trap_vel_comp.png}
    \caption{En este caso, con un Tiempo de Rampa de $t=0.1s$, puede verse como la corriente pico apenas supera los $0.4 A$ }
\end{figure}
El perfil de velocidad se desarrolló implementando un generador de trayectorias, es decir que se permite ingresar una \textit{matriz de consigna de posición} y luego mediante un bloque MATLAB Function se vectoriza el vector de tiempo y de puntos en el espacio, mas detalles en el Anexo ~\ref{anexo:gen_trayectorias} .
Como era de esperarse, aunque la implementacion del perfil de velocidad trapezoidal corrigiera bastante el problema, las discontinuidades en la corriente, y por lo tanto en la aceleración producían Jerk \textit{infinito}, lo que en una articulación real seria inaceptable, por lo que se expandio el Generador de Trayectorias a la generador de un perfil trapezoidal de la aceleración, obteniendo el siguiente resultado.
\begin{figure}[H]
    \centering
    \includegraphics[width=1\linewidth]{imagenes/perfil_trap_accel_comp.png}
    \caption{Elminación de los picos no derivables de la corriente, o aceleración}
\end{figure}

\section*{Mejora del observador mecanico ante perturbaciones constantes}

El observador mecanico originalmente implementado para el punto 5.2.3 responde a la estructura clasica de Luenberger para el subsistema mecanico simplificado, utilizando como variable medida la posicion angular del eje del motor $\theta_m(t)$ y reconstruyendo la velocidad angular $\omega_m(t)$ a partir del modelo y del error de salida. Dado que la consigna establece ubicar los polos del observador en
\[
p_{obs\,1,2}=-3200\ \text{rad/s},
\]
se adopto inicialmente el observador de segundo orden como diseño base.

Partiendo del modelo mecanico simplificado utilizado para el diseño,
\begin{equation}
\dot{\theta}_m(t)=\omega_m(t), 
\qquad
\dot{\omega}_m(t)=\frac{1}{J_{eq}}\,T_m'(t),
\end{equation}
y definiendo la salida medida como
\begin{equation}
y(t)=\theta_m(t),
\end{equation}
el observador base queda dado por:
\begin{equation}
\label{eq:obs_base}
\begin{aligned}
    \dot{\hat{\theta}}_m(t) &= \hat{\omega}_m(t) + K_{e\theta}\big(\theta_m(t)-\hat{\theta}_m(t)\big), \\
    \dot{\hat{\omega}}_m(t) &= \frac{1}{J_{eq}}\,T_m'(t) + K_{e\omega}\big(\theta_m(t)-\hat{\theta}_m(t)\big).
\end{aligned}
\end{equation}

Definiendo el error de estimacion como
\begin{equation}
e_\theta(t)=\theta_m(t)-\hat{\theta}_m(t),
\qquad
e_\omega(t)=\omega_m(t)-\hat{\omega}_m(t),
\end{equation}
la dinamica del error para el observador base queda:
\begin{equation}
\begin{bmatrix}
\dot e_\theta(t)\\
\dot e_\omega(t)
\end{bmatrix}
=
\begin{bmatrix}
-K_{e\theta} & 1\\
-K_{e\omega} & 0
\end{bmatrix}
\begin{bmatrix}
e_\theta(t)\\
e_\omega(t)
\end{bmatrix}.
\end{equation}

Por lo tanto, el polinomio caracteristico asociado es
\begin{equation}
p_{obs}(s)=s^2+K_{e\theta}s+K_{e\omega}.
\end{equation}

Imponiendo como polinomio deseado
\begin{equation}
p_{des}(s)=(s+3200)^2=s^2+6400\,s+3200^2,
\end{equation}
se obtienen las ganancias del observador base:
\begin{equation}
\label{eq:ganancias_obs_base}
K_{e\theta}=6400,
\qquad
K_{e\omega}=3200^2=1.024\times10^7.
\end{equation}

Con estas ganancias, el observador reconstruye adecuadamente la velocidad en transitorios y presenta una rapidez de convergencia suficientemente mayor que la dinamica dominante del controlador externo. No obstante, al verificar el comportamiento exigido en el punto 5.2.5.b, se observa que ante perturbaciones de carga cuasi-estacionarias aparece un error de estimacion de regimen permanente distinto de cero. Esto indica que, aunque el observador base es estable y rapido, no rechaza completamente perturbaciones constantes o pequenos sesgos de modelado.

\begin{figure}[H]
    \centering
    \includegraphics[width=0.78\linewidth]{imagenes/5_2_5_B_error_observador.png}
    \caption{Error de offset observado con el observador mecanico de segundo orden sin accion integral, ante perturbaciones de carga cuasi-estacionarias.}
    \label{fig:obs_offset_original}
\end{figure}

En consecuencia, el observador de segundo orden se conserva como paso intermedio de diseño correspondiente al punto 5.2.3, pero no como solucion definitiva para el punto 5.2.5.b. Para eliminar este offset se propone aumentar el observador con una accion integral del error de posicion.

Se define un nuevo estado auxiliar $z(t)$ como:
\begin{equation}
\label{eq:def_z_obs}
e_{\theta}(t)=\theta_m(t)-\hat{\theta}_m(t),
\qquad
\dot z(t)=e_{\theta}(t).
\end{equation}

De esta manera, el observador mecanico aumentado queda:
\begin{equation}
\label{eq:obs_integral}
\begin{aligned}
    \dot{\hat{\theta}}_m(t) &= \hat{\omega}_m(t) + K_{e\theta}\,e_{\theta}(t), \\
    \dot{\hat{\omega}}_m(t) &= \frac{1}{J_{eq}}\,T_m'(t) + K_{e\omega}\,e_{\theta}(t) + K_{ei}\,z(t), \\
    \dot z(t) &= e_{\theta}(t).
\end{aligned}
\end{equation}

En esta formulacion, el termino $K_{ei}z(t)$ introduce una correccion integral sobre el error de posicion. Fisicamente, esto permite compensar perturbaciones mecanicas aproximadamente constantes o incertidumbres de modelado que, en el observador base, se manifestaban como un sesgo persistente en la velocidad estimada y, por integracion, tambien en la posicion estimada.

Definiendo ahora el vector de error aumentado
\begin{equation}
\mathbf e_a(t)=
\begin{bmatrix}
e_\theta(t)\\
e_\omega(t)\\
z(t)
\end{bmatrix},
\end{equation}
la dinamica del error queda gobernada por
\begin{equation}
\dot{\mathbf e}_a(t)=
\underbrace{
\begin{bmatrix}
-K_{e\theta} & 1 & 0\\
-K_{e\omega} & 0 & -K_{ei}\\
1 & 0 & 0
\end{bmatrix}
}_{\mathbf A_{obs,a}}
\mathbf e_a(t),
\end{equation}
cuyo polinomio caracteristico resulta
\begin{equation}
\label{eq:polinomio_obs_integral}
p_{obs,a}(s)=s^3+K_{e\theta}s^2+K_{e\omega}s+K_{ei}.
\end{equation}

Para conservar aproximadamente la misma rapidez principal de convergencia que en el observador original, se ubican los tres polos reales e iguales en:
\begin{equation}
p_{1,2,3}=-3200\ \text{rad/s}.
\end{equation}

Asi, el polinomio deseado es
\begin{equation}
p_{des,a}(s)=(s+3200)^3
=s^3+9600\,s^2+3.072\times10^7\,s+3.2768\times10^{10}.
\end{equation}

Comparando coeficientes con \eqref{eq:polinomio_obs_integral}, se obtienen las ganancias del observador aumentado:
\begin{equation}
\label{eq:ganancias_obs_integral}
K_{e\theta}=9600,
\qquad
K_{e\omega}=3.072\times10^7,
\qquad
K_{ei}=3.2768\times10^{10}.
\end{equation}

Desde el punto de vista de implementacion en Simulink, esta mejora se incorpora mediante una rama adicional en paralelo:
\begin{enumerate}
    \item se calcula el error de posicion $e_{\theta}(t)=\theta_m(t)-\hat{\theta}_m(t)$;
    \item dicho error alimenta un integrador $1/s$, cuya salida es el estado $z(t)$;
    \item la salida $z(t)$ se multiplica por la ganancia $K_{ei}$;
    \item la senal $K_{ei}z(t)$ se suma en la ecuacion de $\dot{\hat{\omega}}_m(t)$ junto con el termino de modelo y la correccion proporcional $K_{e\omega}e_{\theta}(t)$;
    \item la velocidad estimada $\hat{\omega}_m(t)$ se integra para obtener $\hat{\theta}_m(t)$, cerrando la realimentacion del observador.
\end{enumerate}

\begin{figure}[H]
    \centering
    \includegraphics[width=0.92\linewidth]{imagenes/5_2_5_B_diagrama_bloques_observador.png}
    \caption{Diagrama de bloques del observador mecanico aumentado con accion integral implementado en Simulink. El error de posicion $e_{\theta}(t)=\theta_m(t)-\hat{\theta}_m(t)$ alimenta tanto las ramas proporcionales como la rama integral, cuya salida corrige la dinamica de $\dot{\hat{\omega}}_m(t)$ para eliminar el error estacionario de estimacion.}
    \label{fig:obs_bloques_integral}
\end{figure}

En la practica, esta estructura puede interpretarse tambien como una estimacion implicita de una perturbacion mecanica constante o de un \textit{bias} equivalente no modelado. Sin embargo, independientemente de esa interpretacion fisica, la conclusion relevante para el presente trabajo es que la incorporacion del estado integral elimina el error estacionario de estimacion frente a perturbaciones constantes de carga.

\begin{figure}[H]
    \centering
    \includegraphics[width=0.78\linewidth]{imagenes/5_2_5_B_observador_sin_offset.png}
    \caption{Respuesta del observador aumentado con accion integral. Se observa la desaparicion del error de offset en regimen permanente.}
    \label{fig:obs_offset_corregido}
\end{figure}

Comparando las Fig.~\ref{fig:obs_offset_original} y \ref{fig:obs_offset_corregido}, se verifica que el agregado del estado integral corrige el error de estimacion de regimen permanente sin perder una rapidez de convergencia compatible con el resto de la arquitectura de control. Por este motivo, el observador mecanico aumentado con accion integral se adopta como solucion final para el punto 5.2.5.b.

\subsection*{Punto 5.2.5.c: comportamiento termico del motor}

Para verificar el comportamiento termico del accionamiento en operacion repetitiva, se aplico de manera ciclica la consigna final de posicion adoptada en el trabajo, compuesta por un trapecio de 15 s de duracion: rampa de subida entre 0 y 5 s, mantenimiento entre 5 y 10 s, y rampa de bajada entre 10 y 15 s. Sobre esta condicion de funcionamiento se considero una temperatura ambiente de
\[
T_{amb}=25^\circ\text{C}
\]
y un torque de perturbacion maximo de
\[
T_{Ld}=5\ \text{Nm}.
\]

La simulacion se ejecuto repitiendo sucesivamente dicho ciclo, de modo de evaluar la evolucion de la temperatura del bobinado del estator \(T_s^\circ(t)\). Dado que cada ciclo tiene una duracion de 15 s, el instante en que se supera la temperatura maxima admisible
\[
T_{s,\max}=115^\circ\text{C}
\]
ocurre a los
\[
t=225.9\ \text{s},
\]
lo que corresponde a aproximadamente 15 ciclos completos y 0.9 s del ciclo numero 16.

\begin{figure}[H]
    \centering
    \includegraphics[width=0.9\linewidth]{imagenes/5_2_5_C_comportamiento_termico.png}
    \caption{Evolucion de la temperatura del bobinado \(T_s^\circ(t)\) bajo operacion repetitiva, considerando \(T_{amb}=25^\circ\text{C}\), \(T_{Ld}=5\,\text{Nm}\) y la consigna trapezoidal de posicion utilizada en el trabajo. Se observa que la temperatura supera el limite admisible \(T_{s,\max}=115^\circ\text{C}\) a los 225.9 s.}
    \label{fig:termico_525c}
\end{figure}

Por lo tanto, para las condiciones analizadas, el sistema no satisface el requerimiento termico de operacion continua repetitiva, ya que la temperatura del bobinado excede el valor maximo permitido antes de alcanzar el equilibrio termico. En consecuencia, esta condicion de trabajo no resulta admisible desde el punto de vista termico.
%-------------
\subsubsection{Evaluación de desempeño utilizando sensores no ideales ($G\neq 1$)}
En esta sección se pretende evaluar el desempeño del sistema asumiendo una función de transferencia para los sensores de modo que se comporten como filtros pasa bajos, para ello en la consigna se detalla un punto de partida para los polos de las funciones de transferencia a aplicar mediante el modelo de SS.

Se reemplaza entonces el modelo ideal de medición $G(s)=1$ por modelos dinámicos equivalentes de filtros pasa bajos con ganancia unitaria, implementados en espacio de estados. Los sensores afectados son las corrientes de fase, la posición angular del rotor y la medición de temperatura del estator, con los parámetros indicados en la consigna :contentReference[oaicite:0]{index=0}.

\subsubsection*{Polos requeridos para cada sensor}

Tomando como referencia la forma estándar de un filtro pasa bajos de primer y segundo orden, los polos requeridos para cada sensor son los que se resumen en la Tabla~\ref{tab:polos_sensores}.

\begin{table}[H]
    \centering
    \caption{Polos requeridos para los modelos de sensores no ideales}
    \label{tab:polos_sensores}
    \begin{tabular}{|c|c|c|c|}
        \hline
        \textbf{Sensor} & \textbf{Modelo} & \textbf{Parámetros} & \textbf{Polos} \\
        \hline
        $i_{as}(t),\,i_{bs}(t),\,i_{cs}(t)$ 
        & LP 2$^\circ$ orden 
        & $\omega_n=6000\ \text{rad/s},\ \xi=1$ 
        & $p_{1,2}=-6000$ \\
        \hline
        $\theta_m(t)$ 
        & LP 2$^\circ$ orden 
        & $\omega_n=2000\ \text{rad/s},\ \xi=1$ 
        & $p_{1,2}=-2000$ \\
        \hline
        $T_s^\circ(t)$ 
        & LP 1$^\circ$ orden 
        & $\tau=20\ \text{s}$ 
        & $p=-\dfrac{1}{\tau}=-0.05$ \\
        \hline
    \end{tabular}
\end{table}

Para los filtros de segundo orden se parte de la función de transferencia estándar con ganancia estática unitaria:
\begin{equation}
    G_{LP,2}(s)=\frac{\omega_n^2}{s^2+2\xi\omega_n s+\omega_n^2}.
    \label{eq:lp2_tf}
\end{equation}

Los polos asociados resultan de resolver la ecuación característica
\begin{equation}
    s^2+2\xi\omega_n s+\omega_n^2=0,
\end{equation}
por lo que
\begin{equation}
    p_{1,2}=-\xi\omega_n\pm\omega_n\sqrt{\xi^2-1}.
\end{equation}

Como en este caso se adopta $\xi=1$, ambos filtros de segundo orden quedan críticamente amortiguados, con polo doble en
\begin{equation}
    p_1=p_2=-\omega_n.
\end{equation}

Por otra parte, para el sensor térmico se emplea un filtro pasa bajos de primer orden con ganancia unitaria:
\begin{equation}
    G_{LP,1}(s)=\frac{1}{\tau s+1},
    \label{eq:lp1_tf}
\end{equation}
cuyo polo se ubica en
\begin{equation}
    p=-\frac{1}{\tau}.
\end{equation}

\subsubsection*{Obtención del modelo en espacio de estados}

A continuación se desarrolla la obtención de las matrices $A$, $B$ y $C$ para cada tipo de filtro, partiendo de sus funciones de transferencia estándar.

\subsubsection*{Filtro pasa bajos de primer orden}

Partiendo de \eqref{eq:lp1_tf}:

\begin{equation}
    \frac{Y(s)}{U(s)}=\frac{1}{\tau s+1}.
\end{equation}
\begin{equation}
    (\tau s+1)Y(s)=U(s).
\end{equation}

Aplicando transformada inversa de Laplace:
\begin{equation}
    \tau \dot{y}(t)+y(t)=u(t).
\end{equation}

Despejando la derivada:
\begin{equation}
    \dot{y}(t)=-\frac{1}{\tau}y(t)+\frac{1}{\tau}u(t).
\end{equation}

Si se elige como variable de estado
\begin{equation}
    x(t)=y(t),
\end{equation}
entonces el modelo en espacio de estados queda:
\begin{equation}
    \dot{x}(t)=Ax(t)+Bu(t), \qquad y(t)=Cx(t),
\end{equation}
con
\begin{equation}
    A=\begin{bmatrix}
        -\dfrac{1}{\tau}
    \end{bmatrix},
    \qquad
    B=\begin{bmatrix}
        \dfrac{1}{\tau}
    \end{bmatrix},
    \qquad
    C=\begin{bmatrix}
        1
    \end{bmatrix}.
\end{equation}

Por lo tanto, para el sensor de temperatura con $\tau=20\ \text{s}$:
\begin{equation}
    A_T=\begin{bmatrix}
        -0.05
    \end{bmatrix},
    \qquad
    B_T=\begin{bmatrix}
        0.05
    \end{bmatrix},
    \qquad
    C_T=\begin{bmatrix}
        1
    \end{bmatrix}.
\end{equation}

\subsubsection*{Filtro pasa bajos de segundo orden}

Partiendo de \eqref{eq:lp2_tf}:
\begin{equation}
    \frac{Y(s)}{U(s)}=\frac{\omega_n^2}{s^2+2\xi\omega_n s+\omega_n^2}.
\end{equation}
\begin{equation}
    (s^2+2\xi\omega_n s+\omega_n^2)Y(s)=\omega_n^2 U(s).
\end{equation}

Aplicando transformada inversa de Laplace:
\begin{equation}
    \ddot{y}(t)+2\xi\omega_n \dot{y}(t)+\omega_n^2 y(t)=\omega_n^2 u(t).
\end{equation}

Ahora se define el vector de estado como
\begin{equation}
    x_1(t)=y(t), \qquad x_2(t)=\dot{y}(t).
\end{equation}

Entonces:
\begin{equation}
    \dot{x}_1(t)=x_2(t),
\end{equation}
\begin{equation}
    \dot{x}_2(t)= -\omega_n^2 x_1(t)-2\xi\omega_n x_2(t)+\omega_n^2 u(t).
\end{equation}

En forma matricial:
\begin{equation}
    \dot{x}(t)=Ax(t)+Bu(t), \qquad y(t)=Cx(t),
\end{equation}
con
\begin{equation}
    A=
    \begin{bmatrix}
        0 & 1 \\
        -\omega_n^2 & -2\xi\omega_n
    \end{bmatrix},
    \qquad
    B=
    \begin{bmatrix}
        0 \\
        \omega_n^2
    \end{bmatrix},
    \qquad
    C=
    \begin{bmatrix}
        1 & 0
    \end{bmatrix}.
\end{equation}

\subsubsection*{Aplicación al sensor de corriente}

Para los sensores de corriente se adopta $\omega_n=6000\ \text{rad/s}$ y $\xi=1$, luego:
\begin{equation}
    A_i=
    \begin{bmatrix}
        0 & 1 \\
        -6000^2 & -2\cdot 1 \cdot 6000
    \end{bmatrix}
    =
    \begin{bmatrix}
        0 & 1 \\
        -3.6\times 10^7 & -1.2\times 10^4
    \end{bmatrix},
\end{equation}
\begin{equation}
    B_i=
    \begin{bmatrix}
        0 \\
        6000^2
    \end{bmatrix}
    =
    \begin{bmatrix}
        0 \\
        3.6\times 10^7
    \end{bmatrix},
    \qquad
    C_i=
    \begin{bmatrix}
        1 & 0
    \end{bmatrix}.
\end{equation}

Este mismo modelo se aplica de manera idéntica a $i_{as}(t)$, $i_{bs}(t)$ e $i_{cs}(t)$.

\subsubsection*{Aplicación al sensor de posición}

Para el sensor de posición angular se adopta $\omega_n=2000\ \text{rad/s}$ y $\xi=1$, por lo que:
\begin{equation}
    A_\theta=
    \begin{bmatrix}
        0 & 1 \\
        -2000^2 & -2\cdot 1 \cdot 2000
    \end{bmatrix}
    =
    \begin{bmatrix}
        0 & 1 \\
        -4.0\times 10^6 & -4.0\times 10^3
    \end{bmatrix},
\end{equation}
\begin{equation}
    B_\theta=
    \begin{bmatrix}
        0 \\
        2000^2
    \end{bmatrix}
    =
    \begin{bmatrix}
        0 \\
        4.0\times 10^6
    \end{bmatrix},
    \qquad
    C_\theta=
    \begin{bmatrix}
        1 & 0
    \end{bmatrix}.
\end{equation}

\subsubsection*{Resumen de modelos en espacio de estados}

Finalmente, los modelos de sensores no ideales a implementar quedan definidos como:

\begin{equation}
    \dot{x}_i=A_i x_i + B_i u_i, \qquad y_i=C_i x_i,
\end{equation}
para los sensores de corriente,

\begin{equation}
    \dot{x}_\theta=A_\theta x_\theta + B_\theta u_\theta, \qquad y_\theta=C_\theta x_\theta,
\end{equation}
para el sensor de posición angular, y

\begin{equation}
    \dot{x}_T=A_T x_T + B_T u_T, \qquad y_T=C_T x_T,
\end{equation}
para el sensor de temperatura.

En todos los casos, la entrada $u(t)$ corresponde a la variable física real a medir, mientras que la salida $y(t)$ representa la señal efectivamente medida por el sensor no ideal. Para evitar transitorios artificiales al inicio de la simulación, las condiciones iniciales de cada filtro deben fijarse de forma consistente con el valor inicial de la magnitud medida.

\subsubsection*{Simulación con sensores no ideales}
\begin{figure}[h]
    \centering
    \includegraphics[width=0.8\textwidth]{imagenes/sensores_no_ideales_1.png}
    \caption{\textit{Arriba Posicion y Velocidad, debajo corriente en eje en cuadratura.}Puede verse el efecto del modelo en SS de los sensores, donde le es imposible seguir la consigna de posición. ($T_{ld} = 5 \ [Nm]$)}
    \label{fig:sens_pos_1}
\end{figure}
Es interesante observar en la Fig \ref{fig:sens_pos_1} como con los polos elegidos para el modelo son completamente irracionales, lo que era de esperarse, ya que si recordamos el diseño del modulador de corriente, se habian elegido $p_{a,b,c}=-5000 \ [rad/s]$, por lo que, modelar los sensores en este caso de corriente, con $p_{sens,abc}=-6000 \ [rad/s]$ no es lo suficientemente mas rapido que la dinámica del sistema, pasa lo mismo con los polos del observador. Es por eso que se propone aumentar estos polos de forma conjunta, es decir, $\omega_n\prime  = N \cdot \omega_n ;\ N=1,2,3...$
Luego de experimentos, se obtienen resultados aceptables con $N=3$ (Fig. \ref{fig:sens_pos_2}). Lo que indicaría que los sensores a utilizar en nuestro accionamiento deberian presentar una respuesta por lo menos 3 veces más rápida que los del modelo que se especificaron inicialmente en la consigna.
 \begin{figure}[H]
    \centering
    \includegraphics[width=0.8\textwidth]{imagenes/sensores_no_ideales_2.png}
    \caption{\textit{Arriba Posicion y Velocidad, debajo corriente en eje en cuadratura.} Puede verse como ya el accionamiento sigue con éxito la consigna de solicitada sin oscilaciones producidas por los retardos de los sensores. ($T_{ld} = 5 \ [Nm]$)}
    \label{fig:sens_pos_2}
\end{figure}
\subsubsection*{Análisis de respuesta simulando no idealidad de inversor trifásico de tensión}
En un accionamiento físico real, el inversor está restringido por las limitaciones dinámicas de los dispositivos semiconductores de conmutación (como los transistores MOSFETs o IGBTs) y por el nivel máximo de energía que puede proveer su fuente de alimentación de corriente continua.
Para evaluar de manera fidedigna cómo estas limitaciones degradan el desempeño del sistema, se reemplaza el bloque ideal por un modelo equivalente estructurado en el Espacio de Estados (SS) que incorpora dos restricciones físicas fundamentales: un ancho de banda limitado y la saturación estricta de la tensión.
Se propone a priori, simular esta limitación como un filtro pasa bajo de segundo orden con amortiguación crítica ($\zeta = 1$)y $\omega_n = 6000 \ [rad/s]$, lo que siguiendo el desarrollo del punto antrerior, las matrices $A_v$ $B_v$ del Espacio de Estados quedarían:
\begin{equation*}
    A_v=
    \begin{bmatrix}
        0 & 1 \\
        -\omega_n^2 & -2\xi\omega_n
    \end{bmatrix},
    \qquad
    B_v=
    \begin{bmatrix}
        0 \\
        \omega_n^2
    \end{bmatrix},
    \qquad
    C_v=
    \begin{bmatrix}
        1 & 0
    \end{bmatrix}.
\end{equation*}
 \begin{figure}[H]
    \centering
    \includegraphics[width=0.8\textwidth]{imagenes/inversor_no_ideal_6000.png}
    \caption{\textit{Arriba Posicion y Velocidad, debajo corriente en eje en cuadratura.} Puede verse como el inversor no es lo suficientemente rapido. ($T_{ld} = 5 \ [Nm]$, $\omega_n = 6000 \ [rad/s]$)}
    \label{fig:inv_trifasico_6000}
\end{figure}
Como era de esperarse, la rapides propuesta inicialmente no es suficiente, ya que es solo un $20\%$ más rápido que el lazo de corriente. Se propone, siguiendo el mismo procedimiento que en el modelado de los sensores, duplicar la rapidez de respuesta del inversor a $\omega_n\prime = 12000 \ [rad/s]$
 \begin{figure}[H]
    \centering
    \includegraphics[width=0.8\textwidth]{imagenes/inversor_no_ideal_12000.png}
    \caption{\textit{Arriba Posicion y Velocidad, debajo corriente en eje en cuadratura.} Finalmente se consigue un rendimiento satisfactorio y cumplimiento con las demás condiciones de operación al mismo tiempo.  ($T_{ld} = 5 \ [Nm]$, $\omega_n = 12000 \ [rad/s]$)}
    \label{fig:inv_trifasico_12000}
\end{figure}
\section{Discretización del Controlador}

Para posibilitar la implementación física del algoritmo de control en un microcontrolador industrial o Procesador Digital de Señales (DSP), es imperativo efectuar la transición del diseño desde el tiempo continuo (dominio de Laplace, $s$) al tiempo discreto (dominio $z$). En el entorno digital, el controlador no evalúa las señales de forma ininterrumpida, sino que adquiere muestras y computa ecuaciones en diferencias algebraicas a intervalos regulares de tiempo.

El primer paso crítico en este diseño es la determinación del período de muestreo $T_s$. Según el criterio práctico de diseño de sistemas de control digital, la frecuencia de muestreo $f_s = 1/T_s$ debe ser al menos 10 veces mayor que el ancho de banda del lazo de control más rápido, garantizando así que no se degrade la estabilidad global del sistema . En la arquitectura desarrollada, el lazo proporcional de corrientes posee la dinámica más veloz con un ancho de banda $BW_{cont} = 796 \ \mathrm{Hz}$ . Por consiguiente, el período de muestreo nominal exigido se calcula como:
\begin{equation}
    T_s \leq \frac{1}{10 \cdot 796} \approx 1.25 \times 10^{-4} \ \mathrm{s}
\end{equation}.

\subsection{Método de Integración de Tustin}
Para llevar a cabo la discretización de los elementos con inercia dinámica, específicamente los bloques integradores ($1/s$), se emplea el método de Tustin, también conocido como aproximación bilineal o regla de integración trapezoidal . Este método destaca por su robustez analítica frente a otros métodos numéricos, dado que mapea todo el semiplano izquierdo del plano complejo $s$ estrictamente dentro del círculo unitario en el plano $z$, lo que asegura que si el sistema en tiempo continuo es estable, su equivalente discreto también lo será .

Físicamente, el método aproxima la integral temporal de una señal calculando el área bajo la curva mediante el trazado de trapecios (líneas rectas) entre dos instantes de muestreo consecutivos . Matemáticamente, la transformación bilineal de Tustin se define sustituyendo el operador continuo de Laplace $s$ por su equivalente discreto dependiente del período de muestreo:
\begin{equation}
    s \approx \frac{2}{T_s} \frac{1 - z^{-1}}{1 + z^{-1}}
\end{equation}.

Al aplicar esta sustitución a la función de transferencia de un integrador puro ($H(s) = 1/s$), se obtiene la función de transferencia en el dominio $z$:
\begin{equation}
    H(z) = \frac{Y(z)}{X(z)} = \frac{T_s}{2} \left( \frac{1 + z^{-1}}{1 - z^{-1}} \right)
\end{equation} .

Aplicando la transformada $Z$ inversa a esta expresión, se obtiene la ley de actualización (o ecuación en diferencias recursiva) que el microprocesador debe ejecutar secuencialmente en cada paso de cálculo $k$:
\begin{equation}
    y[k] = y[k-1] + \frac{T_s}{2} \big( x[k] + x[k-1] \big)
\end{equation} .
Donde $x[k]$ y $x[k-1]$ representan la entrada muestreada actual y pasada al integrador, mientras que $y[k]$ y $y[k-1]$ representan la salida (área acumulada) actual y pasada, respectivamente.

\subsubsection{Implementación en la Arquitectura de Control}
En el diagrama general del sistema, el algoritmo de Tustin detallado se aplica exclusivamente para sustituir los bloques con memoria, es decir, la acción integral del controlador PID de movimiento ($K_{sia}/s$) y los múltiples integradores internos del Observador de Estado de Orden Reducido. 

Por otro lado, los bloques netamente algebraicos como el Modulador de Torque y el Modulador de Corrientes carecen de dinámica integradora, por lo que su discretización consiste en operar sus ecuaciones trigonométricas y proporcionales sincronizadas bajo el mismo reloj base $T_s$. Finalmente, la frontera entre el controlador digital y el mundo físico continuo (el motor y los actuadores) se modela empleando Retenedores de Orden Cero (ZOH), los cuales emulan a los conversores digital-analógicos manteniendo retenida la señal de control en forma de escalón durante todo el ciclo $T_s$ .
\begin{figure}[H]
    \centering
    \includegraphics[width=0.8\textwidth]{imagenes/observador_discretizado.png}
    \caption{Diagrama de bloques del observador discretizado. (Incorporacion de ZOH, y reemplazo de integradores continuos por Integradores de Tustin)}
\end{figure}
\begin{figure}[H]
    \centering
    \includegraphics[width=0.8\textwidth]{imagenes/controladorPID_discretizado.png}
    \caption{Diagrama de bloques del controlador de posición discretizado. (Incorporacion de ZOH, y reemplazo de integradores continuos por Integradores de Tustin)}
\end{figure}
\section{Codificación del Software de Control}

Con el fin de validar la viabilidad de la implementación física del sistema diseñado, se procedió a reemplazar el diagrama de bloques del controlador discreto por código programado en lenguaje de texto estructurado (lenguaje MATLAB) encapsulado dentro de un bloque \textit{MATLAB Function} [1]. Este procedimiento emula rigurosamente la programación de bajo nivel que se llevaría a cabo en un microcontrolador industrial o Procesador Digital de Señales (DSP) en una aplicación mecatrónica real [1].

El desarrollo del algoritmo de control se estructuró acatando los siguientes lineamientos fundamentales para sistemas embebidos de tiempo real:

\begin{itemize}
    \item \textbf{Precálculo de variables y optimización computacional:} Para garantizar que el procesador pueda ejecutar la totalidad del lazo de control dentro del estricto margen de tiempo impuesto por el período de muestreo $T_s$, es mandatorio realizar el precálculo de todas las variables que no dependen del instante de ejecución actual [1]. Esto abarca constantes como los coeficientes del integrador de Tustin ($T_s/2$) y la inicialización de la cinemática del generador de trayectorias. A nivel de software, esto se implementó declarando variables de memoria estática (tipo \texttt{persistent}), las cuales se calculan y almacenan exclusivamente durante el primer ciclo de reloj ($t=0$), evitando sobrecargar la unidad aritmético-lógica en cada paso de integración.
    \item \textbf{Arquitectura modular y funciones escalares:} El código fuente fue documentado y organizado empleando funciones escalares desagregadas [1]. Esto permite que el flujo de ejecución sea claro, secuencial y determinístico. En cada interrupción del reloj discreto, el algoritmo ejecuta de forma secuencial: 1) La actualización del Planificador de Trayectorias, 2) La estimación a través del Observador de Estado Reducido, 3) El cálculo de la ley de control de movimiento PID, 4) La resolución del Modulador de Torque y corrientes, y finalmente 5) Las Transformaciones de Park inversas para sintetizar y exportar las consignas de tensión trifásica hacia el retentor de orden cero (ZOH).
    \item \textbf{Validación de desempeño:} Una vez finalizada la codificación, se procedió a ensayar el bloque estructurado en lazo cerrado contra el modelo de la planta física continua. Para validar matemáticamente su correcto funcionamiento, se comparó su comportamiento dinámico frente al del modelo en bloques discretos obtenido en la etapa previa [1]. Las curvas de respuesta de posición, velocidad, torque electromagnético y corrientes de fase de ambos enfoques deben presentar un solapamiento exacto, garantizando así la ausencia de retardos temporales, errores de cuantización introducidos por el código o discrepancias algebraicas.
\end{itemize}
\bibliographystyle{ieeetr}
\bibliography{references}





%%%%%%%% APPENDIX %%%%%%%%



\begin{appendices}
\section{Generador de Trayectorias paramétricamente suavizadas mediante interpolación lineal}
\subsection{Descripción General}
Para el control de posición del motor sincrónico de imanes permanentes (PMSM), se incorporó un \textbf{Generador de Trayectorias} en el esquema de control. Este bloque tiene como objetivo principal suavizar las transiciones entre los distintos puntos de consigna espaciales, generando referencias continuas y derivables de posición ($\theta_{ref}$), velocidad ($\omega_{ref}$) y aceleración ($\alpha_{ref}$) en función del tiempo. 

El algoritmo implementa un perfil de velocidad trapezoidal suavizado, comúnmente denominado \textit{curva en S de 7 tramos}. Esta técnica permite limitar el \textit{jerk} (la derivada temporal de la aceleración), lo cual es fundamental para evitar excitaciones bruscas en la mecánica del rotor y la carga, prevenir picos de corriente excesivos en las fases del estator y mitigar la saturación prematura de los actuadores y lazos de control PI.

\subsection{Dinámica del Perfil de 7 Tramos}
Para trasladar el rotor desde una posición inicial a una final en un intervalo de tiempo dado, el algoritmo divide internamente el movimiento en tres fases principales, que a su vez componen los 7 subtramos característicos de la curva en S:
\begin{enumerate}
    \item \textbf{Fase de Aceleración:} El \textit{jerk} toma un valor constante positivo haciendo que la aceleración suba linealmente hasta su límite. Luego, la aceleración se mantiene constante, y finalmente el \textit{jerk} se vuelve negativo para llevar la aceleración suavemente a cero conforme se alcanza la velocidad máxima de crucero.
    \item \textbf{Fase de Velocidad Constante:} La aceleración y el \textit{jerk} son nulos. El sistema se desplaza a la velocidad máxima calculada para el tramo hasta que es necesario comenzar a frenar.
    \item \textbf{Fase de Desaceleración (Frenado):} De forma simétrica a la primera fase, se aplican perfiles de \textit{jerk} controlados para generar una aceleración negativa suave, frenando el motor hasta alcanzar exactamente la posición de destino con velocidad y aceleración finales nulas.
\end{enumerate}

\subsection{Parámetros y Variables de Configuración}
El bloque procesa los siguientes parámetros para generar las señales en tiempo de simulación:
\begin{itemize}
    \item $t$: Tiempo actual de simulación [s].
    \item \texttt{consigna\_posicion}: Matriz de $N \times 2$ donde la primera columna contiene los instantes de tiempo (nodos temporales) y la segunda indica las posiciones angulares objetivo para cada nodo.
    \item $t_{r1}$ (\texttt{t\_rampa\_vel}): Tiempo total asignado a la fase completa de aceleración (o desaceleración).
    \item $t_{r2}$ (\texttt{t\_rampa\_acel}): Fracción de tiempo dedicada exclusivamente al incremento/decremento lineal de la aceleración (fase de \textit{jerk} distinto de cero).
\end{itemize}

A nivel algorítmico y para optimizar la carga computacional en cada paso de integración de la simulación, el sistema precalcula las distancias relativas entre nodos ($\Delta q$), las duraciones de los tramos ($\Delta t$), la velocidad máxima requerida ($h_1$) y la aceleración máxima ($h_2$) una única vez durante el primer ciclo de ejecución, operando luego mediante la evaluación temporal sobre estos límites preestablecidos.
\label{anexo:gen_trayectorias}
\end{appendices}
\end{document}
